% The appendix, now part of main.tex and printed after the references so that reviewers
% get one PDF and so that every reference between it and the body is a live \ref instead
% of a hard-coded number. It was a second standalone document until 28 Oct; the two-way
% hard-coded numbering that cost is recorded in notes/17 and was the reason for
% paper/check_xrefs.py.
%
% The 2027 call allows an appendix either appended to the main paper or uploaded
% separately, with no page limit on it; the ten-page limit is on the body, and the
% bibliography is excluded from that count.
%
% The data tables are GENERATED: `python -m pilot2.appendix_tables` writes paper/tables/a*.tex
% from results/figure_data.csv and the reports beside it. Do not edit those by hand, and
% regenerate after any re-measurement, so an appendix cell cannot disagree with the main
% text by transcription.
\noindent
The main text refers to the appendix using the section numbers below. \texttt{pilot2/appendix\_tables.py} writes every table marked \emph{generated}
directly from the published measurements rather than transcribing them.
Numbers from the earlier measurement campaign, whose analyses \S\ref{sec:basis} lists one by one,
have not all been re-measured on the current numerical basis. Main-text comparisons
using those controls retain the campaign and placement qualifications stated below.

\section{What was measured when, and on what}
\label{sec:basis}
The host's floating-point behavior changed after a reboot on 18~October. Re-masking
is numerically sensitive (\S\ref{sec:repro}); we therefore report three measurement campaigns separately:
\begin{itemize}\itemsep2pt
\item \textbf{Common-basis remeasurement:} nineteen chromosome-22 cells on 23~October
and seven \emph{E.~coli} cells on 19~October, with 50 records per tier. The controlled
tokenizer sweep was also completed after the reboot as one campaign.
\item \textbf{Membership scoring:} each run's saved scoring passes. In the paired
numerical comparison, rank-1 rates reproduced to the reported precision; this does
not estimate their numerical variance.
\item \textbf{Earlier controls:} the 8--16~October objective grid, placement,
edge-conditioned and downstream audits, and the 28~September capacity/parameterization
controls. These were not rerun on the current basis. Main-text comparisons using them
retain their campaign and placement qualifications; fixed-offset rates are not directly
comparable with drawn-offset rates.
\end{itemize}

\section{Probe geometry and canary construction}

\paragraph{Geometry.}
A canary is 96\nt. The first 48\nt{} form the nominal prefix and the last 48\nt{} are the target.
The BPE prefix protocol uses the preceding complete-token boundary, as detailed in \S\ref{app:measurement}.
Exact-recovery rates count records with all 48 target nucleotides correct. The canary is planted inside a
longer host window of 288\nt{}, everywhere except the tokenizer sweep, which uses 300\nt{} because
300 is divisible by every $k$ from 1 to 6, at an offset drawn independently for each copy.
Repetition tiers are 1, 4 and 16 copies; tier~0 is a never-trained canary of identical construction,
scored in the same run by the same decoder to supply the never-trained control.

\paragraph{Why offsets are sampled rather than fixed.}
With every copy at the same offset, a model with learned absolute positions can score a record from
its position, and coarse tokenizers align to that offset. For a comparable analysis of the three tokenizers shared by both arms, the same variance
decomposition moves the tokenizer's share
from \NUM{15.4\%} under drawn placement to \NUM{59.5\%} under fixed placement, with the
objective's share moving the other way, \NUM{57.7\%} to \NUM{25.7\%} (\S\ref{sec:placement}). Fixed
placement manufactures a tokenizer effect. Every number in the main text uses drawn offsets, with
one exception that is stated wherever it is read: the capacity and parameterization controls of
Table~\ref{tab:capacity}, and the spaced-tokenizer rows inside it, were run on 28 September,
nine days before the offset became an option, and are at fixed placement throughout. These results support within-table comparisons under a common placement condition. Their absolute
extraction rates are not directly comparable to those in the drawn-offset tables, and the character model reaching 100\% at
$2.20$\,\bits{} there is the placement effect, not a counterexample to \S\ref{sec:slots}.

\paragraph{Human-variant-derived canaries.}
For the human arm a canary is a chromosome-22 reference window with rare alternate alleles
substituted, anchored on a variant carried by at most five of the 2{,}548 samples of the
GRCh38 call set. The positional builder does not guarantee that neighboring alleles occur
on the anchor carrier's chromosome; the audit below distinguishes haplotypes and composites.
Ranking candidates use the same construction, while fresh hosts are independently sampled
reference windows. Matched candidate construction helps distinguish training membership
from a generic preference for genomic over uniform-random sequences. Allele frequency is computed from AC/AN because the INFO AF field is rounded
to two decimals; the carrier-count filter uses sample genotypes, not that allele-frequency ratio.
Table~\ref{tab:census} is the set as built. One row is particularly relevant: 248 of the 400
segments are anchored on a variant carried by exactly \emph{one} of the 2{,}548 samples. ``At
most five carriers'' is the filter; a clear majority of these windows are built around a variant
unique to one person in this panel. Linkage would additionally require a recovered sequence
compatible with an observed haplotype and an external identifying panel. Anchor uniqueness
alone does not establish that a composite sequence identifies a person; we measure no
such linkage.

\begin{table}[htbp]
\centering\small
\caption{\emph{Generated.} The human canary set, from the metadata stored beside the sequences in
\texttt{canaries\_chr22.npz}. \emph{Anchor} is the variant the window is centered on; segments also
contain whatever other rare variants fall inside the 96\nt.}
\label{tab:census}
% Generated by pilot2/appendix_tables.py -- do not edit. A8: the human canary set, from the metadata in canaries_chr22.npz.

\begin{tabular}{ll}
\toprule
segments & 400 \\
length & 96\nt, of which the last 48 are the target \\
distinct carriers & 366 individuals of 2{,}548; no carrier contributes more than 4 segments \\
filter & allele frequency $\le$ 0.001, at most 5 carriers \\
rare variants per segment & 2.07 on average, 1--8, 827 in total \\
anchor carried by 1 person & 248 segments \\
\quad by 2 / 3 / 4 / 5 & 64 / 40 / 28 / 20 segments \\
\midrule
anchor allele frequency & 0-0.0002: 247, 0.0002-0.0004: 65, 0.0004-0.0006: 39, 0.0006-0.0008: 28, 0.0008-0.001: 21 \\
segments by variant count & 1: 163, 2: 125, 3: 60, 4: 35, 5: 12, 6: 2, 7: 2, 8: 1 \\
\bottomrule
\end{tabular}

\end{table}

\section{Which planted windows match a chromosome in the cohort}
\label{app:haplo}

\paragraph{Why this needed auditing.}
The builder takes the 96\nt{} window centered on a rare anchor variant and substitutes every rare
alternate allele that falls inside it, from a position-keyed table of rare variants shared across
all anchors. The builder does not check whether the anchor's carrier also carries the neighboring alleles or
which of that individual's two chromosomes carries each allele. With 619{,}639 rare biallelic SNVs on
chromosome~22 the density is about 1.2 per 96\nt{}, so most windows pick up a neighbor: 237 of the
400 carry more than one rare allele (Table~\ref{tab:census}). A window assembled that way can be a
sequence that no chromosome in the cohort has.

\paragraph{What the phased genotypes say.}
For every variant inside a planted window (not only the rare ones), we recorded the exact set
of (sample, chromosome copy) pairs carrying the alternate allele. A window is then the sequence a
given chromosome reads only if that chromosome carries every substituted allele \emph{and} reads
reference at every other variant site inside the window. The first version of this audit tested
only the first half of that condition and reported 164 real haplotypes; requiring the second half
removes 19 of them, because the builder leaves the window's common variant sites at the reference
base and a carrier of one of those reads the window differently. The corrected counts: of the 400
windows, \textbf{145 are chromosome-compatible at the audited sites}, 144 of them single-variant windows plus one composite.
Among the 145 the number of carrying chromosomes is 1 for 86 windows, 2 for 23, 3 for 18, 4 for 10
and 5 for 8; never more than five, so ``carried by at most five'' is exact for this subset.
Twenty more are real \emph{genotypes}, whose alleles all belong to one person but sit on both of
their chromosomes, and \textbf{235 are carried by no chromosome in the cohort}. These are constructed variant combinations, not observed chromosome sequences.

One ambiguity remains: 69 indel or multi-allelic sites fall inside 62 of the 400 windows, and neither pass
can say what a chromosome reads there. Restricted to the 338 windows with none, the classification
is 124 haplotypes, 20 genotypes and 194 chimeras, so the haplotype share changes only from 36.25\% (145/400) to 36.69\% (124/338).

\paragraph{What it changes and what it does not.}
The record-membership question asks whether the model has seen a particular string,
not whose string it is; that question is well posed for all 400 records. The classification defines the interpretations supported by the results. We therefore split
every rate by subset (\path{pilot2/analyze_canary_status.py}) and report the stratification wherever it changes a result. One result changes: single-copy identification in a fresh window, $5/21$, $5/21$ and
$8/21$ per seed on the audited haplotypes against $1/23$, $3/23$ and $0/23$ on windows carried by no chromosome in the cohort (\S\ref{sec:ident}). Identification in the record's own window, verbatim recovery, and the
never-trained controls do not separate.

Two qualifications accompany this stratification in \S\ref{sec:ident} and also apply here. It is
nearly collinear with the number of substituted alleles: every chimera is composite and 144 of
145 haplotypes are single-variant. The allele-count split instead contains 26 one-allele and
24 multi-allele records: across their three repeated seed evaluations, rank-1 counts are
$27/78$ and $6/72$ ($0.35$ and $0.08$). These are descriptive averages over repeated records,
not independent-trial sample sizes; this design cannot attribute the contrast to one description rather than the other. And
the subsets hold 21 and 23 \emph{records}, not 63 and 69 scorings: all three training seeds share
\texttt{data\_seed = 1234} and therefore the same member/control/pool split, so the three seeds
re-score the same records. Rates in \S\ref{sec:ident} are therefore given one seed at a time, where
each record contributes one evaluation and the corresponding rate has a binomial interval. The generated tables keep
the pooled rate with a Wilson interval taken on the record count, which accompanies the pooled rate but is not a confidence interval for that pooled estimate. This
design does not estimate the across-seed component of uncertainty.

\paragraph{The construction we would use again.}
\texttt{pilot2/human\_canary.py --phased} chooses a carrier chromosome first and writes the window
as that chromosome reads it. Doing this correctly needs a second pass over the VCF, because a
window also holds variants too common to pass the rare filter, and leaving those sites at the reference base can produce a sequence absent from the cohort. The phased build uses a separately generated set of 400 window coordinates, only six of which overlap the original set. On the 38{,}400 positions inside these
windows the full carrier sets fit in memory, and the phased set is 400 windows that are each one
real chromosome's own sequence, 252 of them unique to a single chromosome among the cohort's
5{,}096 and none shared by more than five. Thirty-seven common alleles are written in across the
400 windows, a step the original construction did not perform. In this phased set, seventy-three
sites inside 67 windows are indels or not biallelic SNVs and are left at the reference base, so
even the rebuilt windows are each one chromosome's sequence at its SNV sites rather than its every
base. The measurements in this paper use the original set, which stays reproducible; the rebuilt
set is what any extension of this work should plant.

\section{The matched control for the own-window identification rate}
\label{app:bh}

\paragraph{Why a matched control is needed.}
An own-window member and a fresh-window non-member differ in both membership and
contextual access. Their contrast alone cannot distinguish record-specific learning
from a generic advantage of memorized flanks. We therefore evaluate never-trained
records in both borrowed training windows and fresh windows.

\paragraph{The control.}
A non-member has no training row, so it borrows one (\texttt{pilot2/data.py},
\texttt{borrowed\_hosts}): each control is assigned a member's training row and that member's own
offset, and the probe region is overwritten by each candidate exactly as it is for a member. The
conditioning is then identical: the same memorized flanks, the same offset, the same hundred
real rare segments in the pool. The only difference is whether the model has seen the 96
nucleotides now sitting between those flanks. The donor is drawn from the \emph{single-copy}
members without replacement, because the single-copy own-window rate is the one the main text
quotes and that is the host condition the control has to match. \texttt{--donor\_tier} selects the
tier instead, so the four- and sixteen-copy own-window rates, which had no control of their own
when this paragraph was first written, have one now: two further scoring passes over the same 18
cells (\texttt{pilot2/queue\_1028.sh}).

An earlier draft incorrectly described this control as ``the harder case because the memorized
record is competing.'' The control does not create such competition. The ranking pass overwrites
the probe region with each candidate before scoring, so the string the model memorized at that
offset is not among the hundred and one candidates and cannot compete with them. The low control rates
argue against a generic training-window advantage as the sole explanation for the member rates.
What it cannot separate is a model whose conditional distribution at that offset is sharp on a
string absent from the candidate list: that would also drive the control to chance. The claim the
own-window column carries is the first of those, and only the first.

\paragraph{What it says.}
We re-scored every cell of Table~\ref{tab:ident} (\texttt{pilot2/rescore.py},
\texttt{pilot2/queue\_1027\_gen.sh}), retraining nothing and overwriting nothing already on
disk. Re-scoring reproduces every member number to the digit and every held-out loss to four decimals.
The borrowed-host control reads $0.00$ to $0.04$ across all 18 cells against members at $0.04$ to
$0.96$ in the same condition, and $0.04 / 0.00 / 0.02$ at the final checkpoint against
$0.94 / 0.92 / 0.96$. The own-window advantage is specific to the trained record in its trained context,
rather than a generic advantage from familiar flanks. This does not separate record learning
from memorization of a record--host association; fresh-window results address a different access condition.

Matched controls were subsequently added for the two remaining tiers using the same procedure. At four copies the control reads
$0.00$--$0.04$ against members at $0.06$ to $1.00$, and at sixteen it reads $0.00$ in every one of
the 18 cells against members at $0.34$ to $1.00$; at the final checkpoint the members are at
$1.00$ in all three seeds under both tiers. The control therefore does not weaken where the record
was planted more often, which is what it has to do to be one: more copies make the host window
more familiar to the model, and a never-trained record placed in that window still ranks at
chance.

We extended the same control to the other three architectures, which previously lacked it. Nucleotide
Transformer was re-scored from saved weights and reads $0.00$ in all three seeds at both sizes,
against members at $0.64 / 0.66 / 0.60$ on the 250M and $0.74 / 0.90 / 0.86$ on the 500M.
HyenaDNA had to be retrained instead, because that script deleted every checkpoint it saved; the
retrain reproduced the original run's member rates ($0.40 / 0.44 / 0.50$), its fresh-window rates
($0.02$ in all three) and its held-out losses ($2.221 / 2.221 / 2.247$), and its borrowed-host
control is $0.00$ in all three seeds.

The overlapping $k$-mer family is the fourth, and it had no own-window control at all. Its human
runs predate \texttt{--drop\_ckpt} and kept no weights either, so like HyenaDNA they were
retrained rather than re-scored: twelve runs, $k=3$ to $6$ over three seeds, with the control in
place and the checkpoints kept this time. The retrain reproduces the originals \emph{exactly}:
every member own-window and fresh-window rate to the digit and every held-out loss to three
decimals, in all twelve. Its own-window control reads $0.00$--$0.04$ against members at
$0.04$--$0.10$ for $k=3$ to $5$ and $0.24$--$0.26$ at $k=6$, so only the coarsest cut of this
family separates from its control by a margin this many records can carry: four of fifty against
one of fifty is not a separation, and twelve of fifty against one is. In the fresh window none of
them establishes a separation under the tested scorer, candidate pool and sample size: members
read $0.02$--$0.04$ and controls $0.00$--$0.04$. This negative assay does not establish
absence of other membership signals. The stronger fresh-window result in \S\ref{sec:ident}
is established here for DNABERT-2, not for every released genomic family. Verbatim recovery in the own window is zero for $k=3$
to $5$ and $0.06$ at $k=6$.

The evaluation comprised thirty-three cells on 27 October, followed on 28 October by 36 additional
scoring passes and twelve retraining runs (\path{pilot2/queue_1028.sh}). All runs completed without failure.

\paragraph{How to count a control.}
Each row of 900 is 50 records re-scored in 18 cells, not 900 trials, and the two must not be
confused: an earlier version of this paper reported ``3 of 900 never-trained controls ranked
first, below chance'' when all three were one record in one seed at three consecutive rounds.
Under the corrected donor draw the 900 own-window scorings contain 11 rank-1 hits from
\emph{four} distinct records and the 900 fresh-window scorings 3 hits from \emph{one}, and the
two sets overlap: probe~175 ranks first at every one of the six checkpoints of seed~2 in the
own-window condition and at three of them in the fresh-window condition. One control record in
fifty that the candidate pool happens to favor, in one seed, is what that looks like; it is not
a membership signal, and a pooled rate over the row would have hidden it inside a small number.
The interpretable figure is the per-cell count, 0 to 2 of 50, and
\texttt{pilot2/count\_control\_hits.py} prints the distinct records behind any total so that a
total is never quoted without one.

\section{The BPE decoder's normalization was a prior over letters}
\label{app:bpenorm}

\paragraph{What the division did.}
The byte-pair decoder scores candidate letter $a$ at nucleotide $j$ by writing $a$ in,
re-tokenizing, masking every token whose span covers $j$, and reading the log-probability of those
tokens' true identities. Candidates re-segment the window differently, so they mask different
numbers of tokens covering different numbers of nucleotides, and our first version divided by the
nucleotides covered to make the numbers comparable. That divisor depends on the candidate, and a
log-probability is negative, so the letter whose tokens happen to span more nucleotides is pulled toward zero, increasing its likelihood of selection.

\paragraph{How large it is.}
We measure the divisor itself on the released DNABERT-2 vocabulary with no model involved, over 1{,}920 positions in real
planted windows (\texttt{pilot2/diag\_bpe\_denominator.py}): the covering span averages 5.3\nt{} and
the spread across the four candidates averages 1.79\nt{}, reaching 11. At the 1{,}435 positions
where the four divisors differ, the largest falls on A 707 times against 359 expected, with C 312,
T 229 and G 187. ``Continues the run to the left'' accounts for only $0.291$ of those, barely above
chance, so this is a property of the vocabulary and the division, not of homopolymer merging.

\paragraph{What it cost.}
Iterated over 48 positions the prior decides the output. Under the divided score, one letter covers
three quarters or more of the decoded string in \textbf{95\% of 16-copy cells, 85\% of single-copy
cells and 75\% of controls}; the decoded Hamming distance is 34.9, 32.1 and 34.6 of 48 against a
uniform-letter baseline of 36. All three groups retain large reconstruction errors. One cell is worth
naming: seed~0, probe~102 is reported as 48 A's at Hamming 32, chosen by the truth-free selection
step over a restart at Hamming 26, while the leave-one-out diagnostic for that same record reads top-1
$1.00$ at $0.001$\,\bits.

\paragraph{The replacement, and what it changes.}
\texttt{nt\_dist\_bpe\_fixed} picks a nucleotide interval around $j$ that is token-boundary aligned
under all four candidates, scores exactly those nucleotides for each, and divides by nothing: the
masked event is ``this interval reads thus'', the conditioning is the identical remainder of the
window, and the four events differ only at $j$. On the same 60 cells the degenerate fraction falls
to \textbf{0\%, 5\% and 0\%}, and the Hamming distance to \textbf{30.1, 30.6 and 35.8}, so members
have lower mean reconstruction error than controls, whose 35.8 errors are near the uniform-letter baseline of 36. \textbf{Exact recovery
is still zero on both sides.} Removing this measured letter-prior artifact does not produce exact
recovery under the tested search. The leave-one-out diagnostic uses additional true target
context and does not establish that whole-record recovery is possible. The comparison isolates
one scoring defect; it does not attribute the remaining failures to search alone.

One limitation remains unresolved: candidates still decompose the
interval into different numbers of masked tokens, and a product of masked-token marginals is not the
joint, so a segmentation with fewer factors remains mildly favored. The old path is kept as
\texttt{--score pernt} so the submitted numbers stay reproducible. The leave-one-out
\emph{bits} figure is not comparable between the two scorers, because one scores a nucleotide and
the other an interval; only the top-1 accuracy and the Hamming distances are.

\section{The four decoders, and the leave-one-out probe}
\label{sec:decoders}
The protocols differ in conditioning and update unit as well as search strategy.
Methods gives the update equations; \S\ref{app:measurement} supplies implementation details.
\begin{description}\itemsep2pt
\item[Prefix-conditioned:] generate one greedy continuation per record and context,
taking the argmax over admissible content tokens and feeding back predictions. There is
no sampling, beam search, refinement or restart. For masked models,
the unrevealed suffix, including the host beyond the target, stays masked.
\item[Token re-masking:] for NT and non-overlapping masked backbones, update covering
tokens given the surrounding host. Leading revealed letters are constrained; trailing
host letters within the final covering token remain latent. Use 128 restarts and at most
six sweeps, then select by truth-free PLL. Float64 accumulation avoids near-tie reordering
from relative summation changes of about $10^{-7}$.
\item[Nucleotide re-masking:] for overlapping $k$-mers, mask all tokens covering a
position and update one nucleotide consistently with its neighbors.
\item[Span-aware re-masking:] for BPE, re-tokenize each letter candidate. The historical
per-nucleotide-divided scorer and common-interval correction are distinguished in
\S\ref{app:bpenorm}; both give zero exact recovery in their reported evaluations.
\end{description}
The last two searches use four passes, with 16 restarts in the chromosome-22 remeasurement
and 32 in earlier runs. LOO supplies every other true target position. Its accuracy and
per-record power summary are conditional-prediction diagnostics in all arms, not recovery
bounds; \S\ref{app:measurement} explains the endpoint and numerical mismatch. Table~\ref{tab:loo}
reports both summaries alongside measured recovery.

\section{Per-seed cells behind the identification table}

All reported DNABERT-2 prefix exactness and Hamming distances use the hidden-48
coordinate crop, implemented by \texttt{pilot2.target48};\newline
\texttt{results/target48\_audit.json} records changed cells and within-cell errors.
Historical raw score files remain unchanged. Early archival CSV rows without
prediction strings are marked \texttt{legacy\_variable} and are not used for the
reported 48\nt{} evaluations; they cannot be recropped from summary counts alone.


Table~\ref{tab:a4} gives every DNABERT-2 chromosome-22 trajectory behind Figure~\ref{fig:ident}, per seed, including sixteen-copy fresh-window identification and the single-copy-matched controls. It consolidates the former main-text summary table with the complete per-seed data. Table~\ref{tab:target48} gives the corresponding cropped Hamming distances and uncertainties. Every row is 50 records per tier:
rounds 15--120 were scored with 10 during the run, which put the single-copy trajectory two
checkpoints late, and were re-scored at 50 for this table.

\begin{table}[p]
\centering
\captionsetup{skip=2pt}
\renewcommand{\arraystretch}{0.93}\small
\caption{\emph{Generated.} Per-seed DNABERT-2 on chromosome~22. \emph{id} is rank-1 against 100
real rare segments, chance $0.0099$; \emph{own} is the window the record was trained in and
\emph{fresh} one the model never saw. The two control columns are the never-trained records in
those same two conditions: \emph{own} is the borrowed-host control of \S\ref{app:bh},
which did not exist before 27 October. Fresh controls are independently drawn;
only own-window contexts are matched pairwise. \emph{verbatim} is all 48 target nucleotides exact under the
prefix-conditioned protocol, which is the extraction the scoring pass records; the span-aware
decoder reads $0\%$ on this model at both tiers (Table~\ref{tab:loo}), so the 16-copy column here
is the protocol the field uses and not ours.}
\label{tab:a4}\label{tab:ident}
% Ten columns, the widest generated table here, and the only wide one that was not scaled: it ran
% 45pt past the right margin. tabcolsep is already 4pt for this document. Scaled like its two
% neighbours instead.
\resizebox{\textwidth}{!}{% Generated by pilot2/appendix_tables.py -- do not edit. A4: per-seed DNABERT-2 on chromosome 22, both host conditions, both controls.

\begin{tabular}{llccccccccc}
\toprule
& & & \multicolumn{2}{c}{16 copies, id} & \multicolumn{2}{c}{1 copy, id} & \multicolumn{2}{c}{control, id} & \multicolumn{2}{c}{verbatim} \\
\cmidrule(lr){4-5}\cmidrule(lr){6-7}\cmidrule(lr){8-9}\cmidrule(lr){10-11}
checkpoint & seed & \bits & own & fresh & own & fresh & own & fresh & 16 & 1 \\
\midrule
round~15 & 0 & 1.266 & 0.40 & 0.40 & 0.04 & 0.04 & 0.02 & 0.00 & 2\% & 0\% \\
 & 1 & 1.261 & 0.34 & 0.30 & 0.04 & 0.04 & 0.00 & 0.00 & 0\% & 0\% \\
 & 2 & 1.263 & 0.38 & 0.34 & 0.04 & 0.04 & 0.02 & 0.00 & 2\% & 0\% \\
\addlinespace[2pt]
round~30 & 0 & 1.299 & 0.94 & 0.78 & 0.04 & 0.04 & 0.02 & 0.00 & 8\% & 0\% \\
 & 1 & 1.292 & 0.96 & 0.84 & 0.06 & 0.04 & 0.00 & 0.00 & 12\% & 0\% \\
 & 2 & 1.296 & 0.98 & 0.80 & 0.08 & 0.04 & 0.02 & 0.02 & 12\% & 0\% \\
\addlinespace[2pt]
round~50 & 0 & 1.357 & 1.00 & 0.96 & 0.26 & 0.06 & 0.00 & 0.00 & 26\% & 0\% \\
 & 1 & 1.353 & 1.00 & 0.96 & 0.22 & 0.06 & 0.00 & 0.00 & 26\% & 0\% \\
 & 2 & 1.352 & 1.00 & 0.94 & 0.22 & 0.04 & 0.02 & 0.02 & 28\% & 0\% \\
\addlinespace[2pt]
round~80 & 0 & 1.454 & 1.00 & 0.94 & 0.46 & 0.08 & 0.00 & 0.00 & 36\% & 0\% \\
 & 1 & 1.456 & 1.00 & 0.94 & 0.42 & 0.12 & 0.00 & 0.00 & 40\% & 0\% \\
 & 2 & 1.457 & 1.00 & 0.98 & 0.50 & 0.12 & 0.02 & 0.02 & 34\% & 0\% \\
\addlinespace[2pt]
round~120 & 0 & 1.599 & 1.00 & 0.92 & 0.80 & 0.14 & 0.00 & 0.00 & 44\% & 0\% \\
 & 1 & 1.593 & 1.00 & 0.98 & 0.78 & 0.22 & 0.02 & 0.00 & 46\% & 0\% \\
 & 2 & 1.607 & 1.00 & 0.96 & 0.74 & 0.24 & 0.02 & 0.00 & 42\% & 0\% \\
\addlinespace[2pt]
final & 0 & 1.726 & 1.00 & 0.96 & 0.94 & 0.18 & 0.04 & 0.00 & 42\% & 0\% \\
 & 1 & 1.728 & 1.00 & 0.96 & 0.92 & 0.22 & 0.00 & 0.00 & 44\% & 0\% \\
 & 2 & 1.734 & 1.00 & 0.98 & 0.96 & 0.26 & 0.02 & 0.00 & 40\% & 0\% \\
\addlinespace[2pt]
\midrule
\multicolumn{11}{l}{\footnotesize 50 records per cell. 13 of the 36 control cells are non-zero, none above $0.04$; \texttt{pilot2/count\_control\_hits.py} gives the distinct records behind them.} \\
\bottomrule
\end{tabular}
}
\par\vspace{3pt}
\normalsize
\caption{Coordinate-cropped Hamming errors out of 48 for single-copy DNABERT-2
records in fresh windows, 50 members and 50 never-trained controls per cell.
SE is the sample standard deviation divided by $\sqrt{50}$; the difference is
control minus member, with SE $\sqrt{s_m^2/50+s_c^2/50}$ for the separately drawn
groups. These are descriptive within-cell uncertainties; repeated checkpoints
and seeds are not pooled as independent trials. Short predictions count missing
letters as errors. Original token-boundary inputs are unchanged.}
\label{tab:target48}
% Generated by pilot2/appendix_tables.py -- do not edit. Coordinate-cropped hidden48 Hamming distances; saved predictions, no new decoding.

\begin{tabular}{llccc}
\toprule
round & seed & member mean $\pm$ SE & control mean $\pm$ SE & difference $\pm$ SE \\
\midrule
15 & 0 & $31.18\pm1.45$ & $33.54\pm1.19$ & $2.36\pm1.88$ \\
 & 1 & $32.92\pm1.40$ & $33.92\pm1.05$ & $1.00\pm1.75$ \\
 & 2 & $30.98\pm1.59$ & $32.02\pm1.36$ & $1.04\pm2.09$ \\
\addlinespace[2pt]
30 & 0 & $32.02\pm1.45$ & $32.30\pm1.40$ & $0.28\pm2.02$ \\
 & 1 & $31.60\pm1.39$ & $32.00\pm1.34$ & $0.40\pm1.93$ \\
 & 2 & $32.18\pm1.39$ & $32.16\pm1.33$ & $-0.02\pm1.92$ \\
\addlinespace[2pt]
50 & 0 & $30.28\pm1.48$ & $32.30\pm1.29$ & $2.02\pm1.96$ \\
 & 1 & $31.16\pm1.38$ & $31.82\pm1.27$ & $0.66\pm1.88$ \\
 & 2 & $30.78\pm1.49$ & $32.40\pm1.37$ & $1.62\pm2.03$ \\
\addlinespace[2pt]
80 & 0 & $29.32\pm1.62$ & $32.34\pm1.37$ & $3.02\pm2.12$ \\
 & 1 & $30.08\pm1.53$ & $30.34\pm1.52$ & $0.26\pm2.16$ \\
 & 2 & $30.22\pm1.41$ & $32.28\pm1.39$ & $2.06\pm1.98$ \\
\addlinespace[2pt]
120 & 0 & $29.98\pm1.44$ & $32.12\pm1.40$ & $2.14\pm2.01$ \\
 & 1 & $31.12\pm1.44$ & $31.30\pm1.33$ & $0.18\pm1.95$ \\
 & 2 & $30.02\pm1.54$ & $32.12\pm1.43$ & $2.10\pm2.10$ \\
\addlinespace[2pt]
final & 0 & $30.16\pm1.47$ & $32.94\pm1.29$ & $2.78\pm1.95$ \\
 & 1 & $29.58\pm1.46$ & $32.06\pm1.25$ & $2.48\pm1.93$ \\
 & 2 & $30.18\pm1.63$ & $31.08\pm1.37$ & $0.90\pm2.13$ \\
\addlinespace[2pt]
\bottomrule
\end{tabular}

\end{table}


\clearpage

\section{The tokenizer sweep, per seed and at both windows}
Table~\ref{tab:a5} links the former main-text means to their per-seed source:
averaging the 300\nt{} rows gives Figure~\ref{fig:slots}. Never-trained controls recover
zero records at every tokenizer snapshot. Nominal slots are $48/k$; the actual
drawn-offset means are given in \S\ref{sec:measure}.

\begin{table}[htbp]
\centering\small
\caption{\emph{Generated.} Extraction at matched held-out loss, records planted $16\times$, per seed,
at both host-window lengths. The main text averages the three seeds of the 300\nt{} block; the
288\nt{} block is the replication. A dash means the run never reaches that loss; \emph{never} in the
last column means extraction never reaches 100\%.}
\label{tab:a5}\label{tab:slots}
% Generated by pilot2/appendix_tables.py -- do not edit. A5: tokenizer sweep at matched held-out loss, per seed, both windows.

\begin{tabular}{llccccccc}
\toprule
window & tokenizer & slots & seed & @2.05 & @2.10 & @2.20 & @2.40 & cost of 100\% \\
\midrule
300\nt & char & 48 & 0 & --- & --- & --- & --- & never \\
 &  & 48 & 1 & --- & --- & --- & --- & never \\
 &  & 48 & 2 & --- & --- & --- & --- & never \\
\addlinespace[1pt]
 & \kmer{2} & 24 & 0 & 0\% & 0\% & 0\% & 0\% & never \\
 &  & 24 & 1 & 0\% & 0\% & 0\% & 0\% & never \\
 &  & 24 & 2 & 0\% & 0\% & 0\% & 0\% & never \\
\addlinespace[1pt]
 & \kmer{3} & 16 & 0 & 0\% & 0\% & 0\% & 2\% & 4.00 \\
 &  & 16 & 1 & 0\% & 0\% & 0\% & 3\% & 4.09 \\
 &  & 16 & 2 & 0\% & 0\% & 0\% & 2\% & 4.10 \\
\addlinespace[1pt]
 & \kmer{4} & 12 & 0 & 0\% & 0\% & 2\% & 12\% & 2.98 \\
 &  & 12 & 1 & 0\% & 0\% & 3\% & 16\% & 3.01 \\
 &  & 12 & 2 & 0\% & 0\% & 1\% & 9\% & 2.98 \\
\addlinespace[1pt]
 & \kmer{5} & 9.6 & 0 & 0\% & 4\% & 23\% & 60\% & 2.62 \\
 &  & 9.6 & 1 & 0\% & 1\% & 25\% & 62\% & 2.64 \\
 &  & 9.6 & 2 & 0\% & 2\% & 24\% & 62\% & 2.62 \\
\addlinespace[1pt]
 & \kmer{6} & 8 & 0 & 2\% & 18\% & 73\% & 93\% & 2.47 \\
 &  & 8 & 1 & 1\% & 16\% & 66\% & 91\% & 2.47 \\
 &  & 8 & 2 & 1\% & 17\% & 63\% & 91\% & 2.46 \\
\addlinespace[1pt]
\midrule
288\nt & char & 48 & 0 & --- & --- & --- & --- & never \\
 &  & 48 & 1 & --- & --- & --- & --- & never \\
 &  & 48 & 2 & --- & --- & --- & --- & never \\
\addlinespace[1pt]
 & \kmer{2} & 24 & 0 & 0\% & 0\% & 0\% & 0\% & never \\
 &  & 24 & 1 & 0\% & 0\% & 0\% & 0\% & never \\
 &  & 24 & 2 & 0\% & 0\% & 0\% & 0\% & never \\
\addlinespace[1pt]
 & \kmer{3} & 16 & 0 & 0\% & 0\% & 0\% & 4\% & 4.12 \\
 &  & 16 & 1 & 0\% & 0\% & 0\% & 3\% & 4.14 \\
 &  & 16 & 2 & 0\% & 0\% & 0\% & 3\% & 4.18 \\
\addlinespace[1pt]
 & \kmer{4} & 12 & 0 & 0\% & 0\% & 3\% & 15\% & 3.01 \\
 &  & 12 & 1 & 0\% & 0\% & 1\% & 9\% & 3.02 \\
 &  & 12 & 2 & 0\% & 0\% & 1\% & 9\% & 3.00 \\
\addlinespace[1pt]
 & \kmer{6} & 8 & 0 & 1\% & 15\% & 66\% & 90\% & 2.48 \\
 &  & 8 & 1 & 2\% & 20\% & 66\% & 92\% & 2.46 \\
 &  & 8 & 2 & 1\% & 19\% & 62\% & 90\% & 2.48 \\
\addlinespace[1pt]
\bottomrule
\end{tabular}

\end{table}

\section{The downstream probe}
\label{sec:downstream}

Table~\ref{tab:a6} reports frozen linear-probe scores and canary recovery at the selected checkpoints.

\begin{table}[htbp]
\centering\small
\caption{\emph{Generated.} Frozen linear-probe MCC on six tasks across two organisms, at the first
checkpoint, at each model's own held-out optimum and at the end of its run, alongside the $16\times$ canary extraction rate over the same interval.
Each summary gives equal weight to distinct (corpus, model, training seed, checkpoint) configurations after merging task panels; each task uses its available seeds.
Coverage varies by task and checkpoint, so these are descriptive summaries, not a fully paired task panel.}
\label{tab:a6}
\resizebox{\textwidth}{!}{% Generated by pilot2/appendix_tables.py -- do not edit. A6: unique-checkpoint downstream aggregates, from results/dnabert_gue_ds/utility_downstream.json.

\begin{tabular}{llccrrrrrr}
\toprule
model & round & held-out & extracted & promoter & core prom. & TF binding & splice & H3 & H3K4me3 \\
\midrule
DNABERT-2 BPE & 0 (first) & 2.688 & 0.0\% & 0.740 & 0.494 & 0.530 & 0.232 & 0.604 & 0.203 \\
 & 3 (best held-out) & 1.943 & 0.0\% & 0.771 & 0.562 & 0.543 & 0.295 & 0.646 & 0.194 \\
 & 50 (last) & 2.589 & 100.0\% & 0.760 & 0.547 & 0.554 & 0.280 & 0.563 & 0.187 \\
\addlinespace[2pt]
DNABERT 3-mer & 0 (first) & 2.840 & 0.0\% & 0.773 & 0.562 & 0.626 & 0.258 & 0.644 & 0.181 \\
 & 30 (best held-out) & 1.916 & 6.0\% & 0.809 & 0.598 & 0.635 & 0.422 & 0.724 & 0.260 \\
 & 50 (last) & 1.948 & 36.7\% & 0.806 & 0.590 & 0.634 & 0.398 & 0.703 & 0.252 \\
\addlinespace[2pt]
DNABERT 4-mer & 0 (first) & 2.659 & 0.0\% & 0.765 & 0.568 & 0.577 & 0.243 & 0.660 & 0.169 \\
 & 15 (best held-out) & 1.920 & 0.0\% & 0.816 & 0.607 & 0.639 & 0.435 & 0.718 & 0.250 \\
 & 50 (last) & 1.959 & 50.8\% & 0.811 & 0.610 & 0.645 & 0.415 & 0.731 & 0.257 \\
\addlinespace[2pt]
DNABERT 5-mer & 0 (first) & 3.086 & 0.0\% & 0.775 & 0.569 & 0.630 & 0.278 & 0.659 & 0.198 \\
 & 15 (best held-out) & 1.928 & 0.0\% & 0.814 & 0.609 & 0.660 & 0.410 & 0.720 & 0.248 \\
 & 50 (last) & 1.974 & 84.2\% & 0.805 & 0.602 & 0.639 & 0.399 & 0.712 & 0.248 \\
\addlinespace[2pt]
DNABERT 6-mer & 0 (first) & 4.581 & 0.0\% & 0.751 & 0.569 & 0.601 & 0.257 & 0.616 & 0.201 \\
 & 15 (best held-out) & 1.916 & 0.0\% & 0.815 & 0.609 & 0.652 & 0.432 & 0.704 & 0.251 \\
 & 50 (last) & 2.027 & 95.0\% & 0.796 & 0.589 & 0.633 & 0.410 & 0.685 & 0.234 \\
\addlinespace[2pt]
HyenaDNA medium & 0 (first) & 1.967 & 0.0\% & 0.764 & 0.568 & 0.626 & 0.277 & 0.654 & 0.223 \\
 & 3 (best held-out) & 1.915 & 0.0\% & 0.788 & 0.587 & 0.664 & 0.279 & 0.679 & 0.246 \\
 & 50 (last) & 2.746 & 82.5\% & 0.780 & 0.570 & 0.638 & 0.213 & 0.644 & 0.228 \\
\addlinespace[2pt]
\bottomrule
\end{tabular}
}
\end{table}

The generated report \path{results/dnabert_gue_ds/utility_downstream.json} records the contributing configuration counts and task coverage.

\paragraph{Separate fine-tuned-head evaluation.}
In addition to the frozen probes in Table~\ref{tab:a6}, we evaluate fine-tuned classification heads
using a scaled-down version of the published recipe (two rounds over up to 20{,}000
labeled sequences, where the published recipe uses three to five rounds, depending on the task,
over the full split). Absolute fine-tuned-head scores are below the corresponding published values;
we interpret changes within each task under this recipe. On
five tasks they land in the published range: promoter \NUM{0.79}--\NUM{0.86}, core promoter
\NUM{0.59}--\NUM{0.66}, TF binding \NUM{0.61}--\NUM{0.68}, splice \NUM{0.60}--\NUM{0.80}, H3
\NUM{0.68}--\NUM{0.74}. Comparing like with like (our DNABERT-2 fine-tuned head at its best
round per task, against the \textsc{DNABERT-2} row of \citet{zhou2024dnabert2} for the same
dataset), we sit below
the published value on all six tasks, by $0.039$ (H3), $0.049$ (promoter), $0.057$ (TF binding),
$0.069$ (H3K4me3), $0.085$ (core promoter) and $0.095$ (splice). The published H3K4me3 figure is
itself only $0.363$, so absolute MCC is low on that task under both recipes; we read it as such.
(Published values re-checked on 31~Oct~2026 against the arXiv HTML of v2: promoter
$0.8677$, core promoter $0.6937$, splice $0.8499$, H3 $0.7827$ and H3K4me3 $0.3627$, all from its
Table~6, and MCC confirmed as the metric in its Table~12. TF binding is $0.7199$, the Table~6
column for human TF dataset~0, which is the dataset we train on; Table~4's $0.7010$ is the mean
over its five human TF datasets and is not the like-for-like comparator. A draft of 27~October
compared against that mean and printed a $0.038$ gap. The comparison is against the plain
\textsc{DNABERT-2} row and not the further-pre-trained one, whose figures differ by up to $0.05$
and in both directions.)

\paragraph{Aggregation and inclusion.}
The \NUM{252} (model, round, task) cells summarize \NUM{476} individual fine-tuned-head
evaluations. All available MCC values, including zeros, enter the task means; no low-MCC
results are excluded. We calculate frozen-probe and fine-tuned-head changes separately,
comparing their available configuration means at each model's selected held-out optimum
and final checkpoint.

\section{The nineteen single-basis chromosome-22 cells}
\label{sec:rerun}

The table below reports every cell in the 23~October batch directly from the measurement reports. \texttt{standard} is the
prefix-conditioned protocol, which the overlapping-$k$-mer and BPE reports do not run;
\texttt{ours} is the decoder matched to that vocabulary. The three tiers of each checkpoint appear
together, and $r{=}0$ is the never-trained control. The \emph{Power diag.} column is the mean
of $p_i^{s_i}$, while \emph{LOO top-1} is the mean of $p_i$. These are conditional-prediction
diagnostics, not recovery bounds (\S\ref{sec:decoders}). Exponentiating the rounded adjacent
accuracy does not reproduce the mean power: the single-copy DNABERT-2 row displays
$0.94$ and $0.268$, while $0.94^{48}=0.05$. Definitions use unrounded per-record values.

\Needspace{125pt}
{\small% Generated by pilot2/appendix_tables.py -- do not edit. A7: the 19 single-basis chromosome-22 cells of 23 Oct.

\begin{longtable}{lcccccc}
\caption{\emph{Generated.} The 23 October re-measurement of chromosome~22: every cell, 50 probes per tier, one host state. \texttt{standard} is the prefix-conditioned protocol, which the overlapping-$k$-mer and BPE reports do not run; \texttt{ours} is the decoder matched to that vocabulary; $r{=}0$ is the never-trained control.}\label{tab:a7}\\
\toprule
run, checkpoint & tier & $n$ & standard & LOO top-1 & power diag. & ours \\
\midrule
\endfirsthead
\multicolumn{7}{l}{\tablename~\thetable\ (continued)}\\
\toprule
run, checkpoint & tier & $n$ & standard & LOO top-1 & power diag. & ours \\
\midrule
\endhead
\bottomrule
\endlastfoot
dnabert2bpe\_human\_s0, ep120 & r=0 & 50 & --- & 0.44 & 0.000 & 0\% \\*
dnabert2bpe\_human\_s0, ep120 & r=1 & 50 & --- & 0.91 & 0.146 & 0\% \\*
dnabert2bpe\_human\_s0, ep120 & r=16 & 50 & --- & 0.99 & 0.882 & 0\% \\
dnabert2bpe\_human\_s0, ep30 & r=0 & 50 & --- & 0.46 & 0.000 & 0\% \\*
dnabert2bpe\_human\_s0, ep30 & r=1 & 50 & --- & 0.61 & 0.001 & 0\% \\*
dnabert2bpe\_human\_s0, ep30 & r=16 & 50 & --- & 0.96 & 0.312 & 0\% \\
dnabert2bpe\_human\_s0, epfinal & r=0 & 50 & --- & 0.41 & 0.000 & 0\% \\*
dnabert2bpe\_human\_s0, epfinal & r=1 & 50 & --- & 0.94 & 0.268 & 0\% \\*
dnabert2bpe\_human\_s0, epfinal & r=16 & 50 & --- & 1.00 & 0.914 & 0\% \\
dnabert2bpe\_human\_s1, ep120 & r=0 & 50 & --- & 0.43 & 0.000 & 0\% \\*
dnabert2bpe\_human\_s1, ep120 & r=1 & 50 & --- & 0.91 & 0.116 & 0\% \\*
dnabert2bpe\_human\_s1, ep120 & r=16 & 50 & --- & 0.99 & 0.866 & 0\% \\
dnabert2bpe\_human\_s1, ep30 & r=0 & 50 & --- & 0.46 & 0.000 & 0\% \\*
dnabert2bpe\_human\_s1, ep30 & r=1 & 50 & --- & 0.61 & 0.002 & 0\% \\*
dnabert2bpe\_human\_s1, ep30 & r=16 & 50 & --- & 0.96 & 0.318 & 0\% \\
dnabert2bpe\_human\_s1, epfinal & r=0 & 50 & --- & 0.43 & 0.000 & 0\% \\*
dnabert2bpe\_human\_s1, epfinal & r=1 & 50 & --- & 0.93 & 0.209 & 0\% \\*
dnabert2bpe\_human\_s1, epfinal & r=16 & 50 & --- & 0.99 & 0.871 & 0\% \\
dnabert2bpe\_human\_s2, ep120 & r=0 & 50 & --- & 0.41 & 0.000 & 0\% \\*
dnabert2bpe\_human\_s2, ep120 & r=1 & 50 & --- & 0.92 & 0.185 & 0\% \\*
dnabert2bpe\_human\_s2, ep120 & r=16 & 50 & --- & 0.99 & 0.801 & 0\% \\
dnabert2bpe\_human\_s2, ep30 & r=0 & 50 & --- & 0.46 & 0.000 & 0\% \\*
dnabert2bpe\_human\_s2, ep30 & r=1 & 50 & --- & 0.60 & 0.001 & 0\% \\*
dnabert2bpe\_human\_s2, ep30 & r=16 & 50 & --- & 0.96 & 0.350 & 0\% \\
dnabert2bpe\_human\_s2, epfinal & r=0 & 50 & --- & 0.42 & 0.000 & 0\% \\*
dnabert2bpe\_human\_s2, epfinal & r=1 & 50 & --- & 0.95 & 0.388 & 0\% \\*
dnabert2bpe\_human\_s2, epfinal & r=16 & 50 & --- & 0.99 & 0.846 & 0\% \\
dnabert3\_nt0.05\_lr0.0001\_human\_s0, ep120 & r=0 & 50 & --- & 0.51 & 0.003 & 0\% \\*
dnabert3\_nt0.05\_lr0.0001\_human\_s0, ep120 & r=1 & 50 & --- & 0.66 & 0.013 & 0\% \\*
dnabert3\_nt0.05\_lr0.0001\_human\_s0, ep120 & r=16 & 50 & --- & 1.00 & 0.873 & 2\% \\
dnabert4\_nt0.05\_lr0.0001\_human\_s0, ep120 & r=0 & 50 & --- & 0.51 & 0.001 & 0\% \\*
dnabert4\_nt0.05\_lr0.0001\_human\_s0, ep120 & r=1 & 50 & --- & 0.63 & 0.011 & 0\% \\*
dnabert4\_nt0.05\_lr0.0001\_human\_s0, ep120 & r=16 & 50 & --- & 0.99 & 0.828 & 0\% \\
dnabert5\_nt0.05\_lr0.0001\_human\_s0, ep120 & r=0 & 50 & --- & 0.50 & 0.000 & 0\% \\*
dnabert5\_nt0.05\_lr0.0001\_human\_s0, ep120 & r=1 & 50 & --- & 0.61 & 0.007 & 0\% \\*
dnabert5\_nt0.05\_lr0.0001\_human\_s0, ep120 & r=16 & 50 & --- & 0.99 & 0.782 & 4\% \\
dnabert6\_nt0.05\_lr0.0001\_human\_s0, ep120 & r=0 & 50 & --- & 0.47 & 0.001 & 0\% \\*
dnabert6\_nt0.05\_lr0.0001\_human\_s0, ep120 & r=1 & 50 & --- & 0.65 & 0.019 & 0\% \\*
dnabert6\_nt0.05\_lr0.0001\_human\_s0, ep120 & r=16 & 50 & --- & 0.99 & 0.704 & 0\% \\
nt250m\_mlm\_human\_s0, ep15 & r=0 & 50 & 0\% & 0.10 & 0.001 & 0\% \\*
nt250m\_mlm\_human\_s0, ep15 & r=1 & 50 & 0\% & 0.18 & 0.005 & 0\% \\*
nt250m\_mlm\_human\_s0, ep15 & r=16 & 50 & 6\% & 0.51 & 0.179 & 6\% \\
nt250m\_mlm\_human\_s0, ep30 & r=0 & 50 & 0\% & 0.09 & 0.000 & 0\% \\*
nt250m\_mlm\_human\_s0, ep30 & r=1 & 50 & 0\% & 0.29 & 0.038 & 0\% \\*
nt250m\_mlm\_human\_s0, ep30 & r=16 & 50 & 10\% & 0.79 & 0.456 & 28\% \\
nt250m\_mlm\_human\_s0, ep50 & r=0 & 50 & 0\% & 0.07 & 0.000 & 0\% \\*
nt250m\_mlm\_human\_s0, ep50 & r=1 & 50 & 2\% & 0.57 & 0.139 & 8\% \\*
nt250m\_mlm\_human\_s0, ep50 & r=16 & 50 & 18\% & 0.96 & 0.809 & 52\% \\
nt500m\_mlm\_human\_s0, ep15 & r=0 & 50 & 0\% & 0.10 & 0.001 & 0\% \\*
nt500m\_mlm\_human\_s0, ep15 & r=1 & 50 & 0\% & 0.24 & 0.013 & 0\% \\*
nt500m\_mlm\_human\_s0, ep15 & r=16 & 50 & 16\% & 0.55 & 0.286 & 14\% \\
nt500m\_mlm\_human\_s0, ep3 & r=0 & 50 & 0\% & 0.08 & 0.000 & 0\% \\*
nt500m\_mlm\_human\_s0, ep3 & r=1 & 50 & 0\% & 0.15 & 0.008 & 0\% \\*
nt500m\_mlm\_human\_s0, ep3 & r=16 & 50 & 2\% & 0.24 & 0.020 & 0\% \\
nt500m\_mlm\_human\_s0, ep50 & r=0 & 50 & 0\% & 0.07 & 0.000 & 0\% \\*
nt500m\_mlm\_human\_s0, ep50 & r=1 & 50 & 0\% & 0.67 & 0.243 & 12\% \\*
nt500m\_mlm\_human\_s0, ep50 & r=16 & 50 & 24\% & 0.96 & 0.838 & 58\% \\
\end{longtable}
}

\section{Numerical reproducibility}
\label{sec:repro}

\paragraph{What moved and what did not.}
Re-scoring a saved checkpoint after the host's floating-point behavior changed
(\texttt{pilot2/rescore.py}) reproduces identification \emph{to the digit} (rank-1 $0.800$ before
and after, held-out loss identical to four decimals) while the extraction rate on the same
checkpoint moves by six points. Small numerical changes can flip near-tied candidate ranks;
a re-masking decode is a discrete iteration in which one argmax resolving differently sends a
restart to a different fixed point. Six points is the same order as the binomial standard error on
50 records, so record-count uncertainty alone does not account for this observed numerical sensitivity.
A single comparison does not estimate its variance or justify a wider numerical interval.
The main text distinguishes the six-point reproducibility observation
from the unchanged identification rates.

\paragraph{Decoder optimization provides a 41$\times$ speedup.}
Batching the 128 restarts into one forward pass per (position, sweep, chunk) and removing about two
million host synchronizations took the audit from 141\,s to 3.4\,s per probe, and the optimized code
reproduces the original bitwise on a full cell (\texttt{--chunk 1}, 8/8 identical). The
re-measurement of nineteen cells took 2\,h43 in total where the original code would have needed
roughly 28 GPU-hours, enabling the rerun on a common numerical basis within the available computation.

\section{Canary placement and position encoding}
\label{sec:placement}

\begin{figure}[ht]
\centering
\includegraphics[width=0.6\textwidth]{figures/F3_variance.pdf}
\caption{Share of the variance in canary memorization by factor over the controlled grid (24 runs).
With independently drawn canary offsets the objective dominates; with every copy at the same offset
the tokenizer appears to, because a model with learned absolute positions can score a record from
its position. Restricted to the three tokenizers both arms share, the drawn-offset shares are
\NUM{57.7}, \NUM{15.4}, \NUM{26.9} and \NUM{0.0}\,\%.}
\label{fig:variance}
\end{figure}

Figure~\ref{fig:variance} reports both decompositions because their difference is larger than most effects that would be
considered substantial within either analysis. This finding concerns the measurement protocol rather
than the models themselves.
Table~\ref{tab:placement} separates the two explanations. If fixed placement inflated memorization
only because the model has learned absolute positions, then a position encoding without them should
remove the effect. It does not: with RoPE, fixed offsets still permit extraction that drawn offsets do not at the same matched loss. The placement effect therefore persists without learned absolute position embeddings.
This control supports the drawn-offset comparison; other
interactions with positional encoding remain possible.

\begin{table}[htbp]
\centering\small
\caption{\emph{Generated.} Extraction of records planted $16\times$ at matched held-out loss, under
absolute and rotary position encodings and under drawn and fixed canary offsets. 20 records per
tier; a dash means the run never reaches that loss.}
\label{tab:placement}
% Generated by pilot2/appendix_tables.py -- do not edit. A9: position encoding against canary placement.

\begin{tabular}{lllcccccc}
\toprule
positions & offsets & tokenizer & seeds & @2.05 & @2.10 & @2.20 & @2.40 & cost of 100\% \\
\midrule
absolute & drawn & char & 2 & 0\% & --- & --- & --- & never \\
 &  & \kmer{3} & 2 & 0\% & 0\% & 0\% & 2\% & 4.12 \\
 &  & \kmer{6} & 2 & 5\% & 35\% & 79\% & 100\% & 2.39 \\
\addlinespace[2pt]
RoPE & drawn & char & 2 & 0\% & --- & --- & --- & never \\
 &  & \kmer{3} & 2 & 0\% & 0\% & 0\% & 2\% & 4.29 \\
 &  & \kmer{6} & 2 & 5\% & 41\% & 84\% & 100\% & 2.40 \\
\addlinespace[2pt]
RoPE & fixed & char & 2 & 0\% & --- & --- & --- & never \\
 &  & \kmer{3} & 2 & 31\% & 61\% & 65\% & 74\% & 3.46 \\
 &  & \kmer{6} & 2 & 100\% & 100\% & 100\% & 100\% & 2.05 \\
\bottomrule
\end{tabular}

\end{table}

\section{Capacity and parameterization controls}
\label{app:capacity}

In the main 300\nt{} non-overlapping tokenizer sweep, all 18 runs record exactly
12,610,560 non-embedding parameters in \texttt{train\_summary.json}.
Total parameter counts for character, 2-, 3-, 4-, 5- and 6-mer tokenizers are
12,768,256; 12,697,600; 12,696,576; 12,782,080; 13,167,616; and 14,735,360,
respectively, identical across the three seeds for each tokenizer. Token and position
embedding tables vary, so this sweep holds the non-embedding backbone fixed rather
than total parameter count. Three controls test whether the tokenizer effect is attributable to
the token embedding table: factorizing the embedding to rank~512 is a pure
reparameterization ($E' = EW^{\!\top}$ with the same tie between input and output, hence the same
function class), yet it moves the \kmer{6} from \NUM{0.61} to \NUM{1.44}\,\bits{} at matched
held-out loss, roughly the \kmer{3} level, and ranks 27 to 256 behave the same, so there is no
dose--response in embedding parameters. It is also not a learning-rate artifact of the embedding
rows: lowering only the embedding learning rate increases measured \kmer{6}
memorization per unit of overfitting in these controls. The character-to-\kmer{3}
difference in the main sweep also occurs without a larger total parameter count.
In the 300\nt{} sweep, they differ by 0.56\%, with the character model slightly larger.

\paragraph{HyenaDNA parameter accounting.}
The released medium-160k model has 6,554,624 trainable parameters, including
its filter biases, matching our training metadata. Its checkpoint stores
14,236,768 values: the remaining 7,682,144 are registered buffers, comprising
7,680,096 positional-grid values (\texttt{pos\_emb.t} and \texttt{pos\_emb.z})
and 2,048 fixed modulation-decay values (\texttt{modulation.deltas}). Counting
all stored values therefore overstates trainable model size; we report 6.55M.

\paragraph{Scaling beyond an extrapolated capacity estimate.}
At $20\times$ data (100{,}000 windows, about 57.6M bits, against roughly 45M bits obtained by
applying the 3.6 bits per parameter of \citet{morris2025capacity} to 12.6M parameters,
an extrapolation from GPT-style models rather than a measurement of ours), exact
recovery and membership scoring separate. At the final checkpoint, the character
model has zero own-window rank-1 and exact-recovery rates at each planted tier in
both seeds. The \kmer{3} also has zero exact recovery, but its sixteen-copy
own-window rank-1 rates are 0\% and 20\% in seeds 0 and 1, respectively.
The \kmer{6} ranks every planted member first in its training window at the final
checkpoint and extracts every sixteen-copy target from round~8 in both seeds.
These measurements are recorded in \path{results/big20x/gate1_report_final.md}
and \path{results/big20x/trajectory.md}. At a fixed round count, $20\times$ data also
means $20\times$ more updates between canary exposures, leaving capacity and
interference confounded. These results do not establish an absence of learning in
the models with zero exact recovery.

Table~\ref{tab:capacity} instead reports the embedding and parameterization controls.
Its experimental conditions require explicit qualification: these runs
predate the placement fix of \S\ref{sec:placement}, so every canary copy sits at a fixed offset, and
they use 20 records per tier rather than 50. Fixed placement substantially increases recovery in these controls, including
saturation across the tested 6-mer variants. The column to read is therefore the last one, the held-out loss at
which extraction first reaches 100\%: every 6-mer variant pays between $2.02$ and $2.08$ whatever
its embedding rank, feed-forward width or embedding learning rate, while the 3-mer pays $2.44$ and
the character model $2.15$. The loss at which full extraction occurs changes with the tokenizer, but not with the
parameterization.

\begin{table}[htbp]
\centering\small
\caption{\emph{Generated.} Capacity and parameterization controls, extraction of records planted
$16\times$ at matched held-out loss, seed means. A dash means the run never reaches that loss.
Fixed canary placement, 20 records per tier; see the caution above.}
\label{tab:capacity}
% Generated by pilot2/appendix_tables.py -- do not edit. A9: capacity and parametrisation controls at matched held-out loss.

\begin{tabular}{lcccccc}
\toprule
variant & seeds & @2.05 & @2.10 & @2.20 & @2.40 & cost of 100\% \\
\midrule
6-mer, unchanged (8 slots) & 1 & 100\% & 100\% & 100\% & 100\% & 2.03 \\
\quad embedding rank 512 & 2 & 100\% & 100\% & 100\% & 100\% & 2.04 \\
\quad rank 256 & 2 & 91\% & 100\% & 100\% & 100\% & 2.06 \\
\quad rank 128 & 2 & 92\% & 100\% & 100\% & 100\% & 2.06 \\
\quad rank 64 & 2 & 84\% & 100\% & --- & --- & 2.08 \\
\quad rank 27 & 2 & 97\% & --- & --- & --- & 2.05 \\
\quad rank 27, 100 epochs & 2 & 92\% & 100\% & 100\% & 100\% & 2.07 \\
\quad feed-forward width 1568 & 2 & 100\% & 100\% & 100\% & 100\% & 2.04 \\
\quad embedding LR $\times 0.3$ & 2 & 100\% & 100\% & --- & --- & 2.02 \\
\quad embedding LR $\times 0.1$ & 2 & 100\% & --- & --- & --- & 2.03 \\
\quad spaced 6-mer (8 slots, 48 positions) & 2 & 87\% & 100\% & 100\% & --- & 2.06 \\
3-mer, embedding rank 512 (16 slots) & 2 & 97\% & 96\% & 95\% & 97\% & 2.44 \\
\quad spaced 3-mer (16 slots, 48 positions) & 2 & 100\% & 100\% & 100\% & 100\% & 2.02 \\
character, embedding rank 512 (48 slots) & 2 & 87\% & 94\% & 100\% & 100\% & 2.15 \\
\bottomrule
\end{tabular}

\end{table}

\section{Edge-conditioned extraction audit}

A masked model does not have to be read left to right, and prefix-only probing is the weak way to
ask it: on diffusion language models, edge-conditioned masks recover up to three times more verbatim
sequences than prefix-conditioned ones \cite{wang2026infilling}. We therefore re-ran every masked
cell of the controlled grid with the sequence after the record left visible and the 48 masked
nucleotides decoded in confidence order. The audit covers five datasets, four tokenizers, both mask rates and two to three seeds,
evaluating never-trained controls and records inserted once or sixteen times
(88 configurations, 264 tier rows, 20 canaries each). Edge-conditioned
extraction recovers \textbf{no} record (Table~\ref{tab:a10}): \NUM{0\%} under both protocols, with
\NUM{34.9} wrong nucleotides out of 48 under the standard protocol and \NUM{35.2} under infilling,
against \NUM{35.4} and \NUM{35.7}, respectively, for non-members. This establishes the from-scratch zero-recovery observation of \S\ref{sec:blind}
under a second tested decoding protocol as well. It does not generalize to all masked models:
the released NT checkpoints yield exact recovery (Table~\ref{tab:human}).
{\small% Generated by pilot2/appendix_tables.py -- do not edit. Grouped edge-conditioned audit; full 264 rows in t_infill_cells.csv.

\begin{longtable}{lllccccccc}
\caption{\emph{Generated.} Edge-conditioned audit across five datasets, four tokenizers and two masking rates. All prefix and infilling exact-recovery rates are zero. Entries are mean Hamming errors out of 48, averaged over the indicated seeds; each seed/tier row evaluates 20 records. $r=0$ denotes never-trained controls. This groups 88 configurations and 264 tier rows, not independent trials. The full rows are retained in \texttt{tables/t\_infill\_cells.csv} and the source report \texttt{results/human\_mlm/infill\_audit.md}.}\label{tab:a10}\\
\toprule
& & & & \multicolumn{2}{c}{$r=0$} & \multicolumn{2}{c}{$r=1$} & \multicolumn{2}{c}{$r=16$} \\
dataset & tokenizer & mask rate & seeds & prefix & infill & prefix & infill & prefix & infill \\
\midrule
\endfirsthead
\toprule
& & & & \multicolumn{2}{c}{$r=0$} & \multicolumn{2}{c}{$r=1$} & \multicolumn{2}{c}{$r=16$} \\
dataset & tokenizer & mask rate & seeds & prefix & infill & prefix & infill & prefix & infill \\
\midrule
\endhead
\bottomrule
\endlastfoot
ecoli & 3mer & 0.15 & 2 & 34.8 & 35.5 & 34.9 & 35.8 & 36.2 & 36.4 \\
ecoli & 3mer & 0.5 & 2 & 35.6 & 36.0 & 35.0 & 35.1 & 35.2 & 36.2 \\
ecoli & 4mer & 0.15 & 2 & 35.7 & 36.4 & 35.1 & 35.6 & 36.4 & 36.0 \\
ecoli & 4mer & 0.5 & 2 & 35.0 & 35.9 & 35.2 & 35.0 & 36.5 & 36.7 \\
ecoli & 6mer & 0.15 & 2 & 35.3 & 35.0 & 35.0 & 36.0 & 35.7 & 36.8 \\
ecoli & 6mer & 0.5 & 2 & 36.0 & 36.2 & 35.7 & 35.1 & 36.5 & 35.7 \\
ecoli & char & 0.15 & 2 & 34.6 & 35.6 & 35.6 & 36.3 & 36.4 & 36.0 \\
ecoli & char & 0.5 & 2 & 36.0 & 35.9 & 35.6 & 35.2 & 36.0 & 36.0 \\
gue & 3mer & 0.15 & 2 & 36.3 & 36.5 & 35.5 & 35.2 & 35.6 & 36.1 \\
gue & 3mer & 0.5 & 2 & 36.4 & 36.0 & 35.5 & 35.5 & 35.2 & 36.2 \\
gue & 4mer & 0.15 & 2 & 36.7 & 36.5 & 34.9 & 35.6 & 35.4 & 36.0 \\
gue & 4mer & 0.5 & 2 & 37.0 & 36.5 & 34.9 & 35.2 & 35.0 & 35.2 \\
gue & 6mer & 0.15 & 2 & 36.0 & 35.9 & 35.3 & 35.7 & 35.2 & 35.4 \\
gue & 6mer & 0.5 & 2 & 37.0 & 36.4 & 34.5 & 35.3 & 34.3 & 35.5 \\
gue & char & 0.15 & 2 & 36.2 & 36.4 & 35.2 & 35.3 & 35.3 & 37.0 \\
gue & char & 0.5 & 2 & 36.2 & 36.4 & 35.5 & 35.2 & 35.3 & 36.8 \\
human & 3mer & 0.15 & 3 & 35.2 & 35.5 & 32.0 & 32.8 & 34.6 & 34.8 \\
human & 3mer & 0.50 & 3 & 32.0 & 34.6 & 31.8 & 33.5 & 33.9 & 33.8 \\
human & 4mer & 0.15 & 3 & 33.7 & 35.3 & 32.0 & 32.7 & 33.1 & 34.2 \\
human & 4mer & 0.50 & 3 & 31.7 & 32.2 & 29.6 & 32.0 & 32.4 & 33.4 \\
human & 6mer & 0.15 & 3 & 35.4 & 35.3 & 31.7 & 32.2 & 34.1 & 34.3 \\
human & 6mer & 0.50 & 3 & 34.4 & 35.3 & 31.9 & 33.1 & 33.3 & 33.9 \\
human & char & 0.15 & 3 & 34.0 & 34.9 & 32.3 & 32.1 & 34.5 & 33.9 \\
human & char & 0.50 & 3 & 34.8 & 35.7 & 31.6 & 32.0 & 34.4 & 34.8 \\
synthetic & 3mer & 0.15 & 2 & 36.1 & 36.0 & 36.4 & 35.8 & 35.2 & 36.0 \\
synthetic & 3mer & 0.5 & 2 & 35.7 & 35.9 & 36.2 & 36.0 & 36.1 & 36.5 \\
synthetic & 4mer & 0.15 & 2 & 35.5 & 36.0 & 36.3 & 36.0 & 36.3 & 35.2 \\
synthetic & 4mer & 0.5 & 2 & 36.8 & 36.5 & 36.0 & 35.0 & 36.0 & 36.1 \\
synthetic & 6mer & 0.15 & 2 & 36.3 & 35.7 & 37.2 & 36.5 & 36.0 & 35.7 \\
synthetic & 6mer & 0.5 & 2 & 35.5 & 36.2 & 36.0 & 36.7 & 35.9 & 36.5 \\
synthetic & char & 0.15 & 2 & 36.3 & 36.4 & 35.1 & 35.2 & 35.1 & 35.0 \\
synthetic & char & 0.5 & 2 & 36.1 & 36.8 & 36.2 & 35.8 & 35.7 & 35.2 \\
yeast & 3mer & 0.15 & 2 & 35.4 & 35.4 & 35.9 & 36.0 & 37.4 & 36.7 \\
yeast & 3mer & 0.5 & 2 & 35.2 & 35.0 & 35.4 & 36.0 & 37.2 & 36.9 \\
yeast & 4mer & 0.15 & 2 & 35.5 & 35.2 & 35.4 & 36.0 & 37.1 & 36.8 \\
yeast & 4mer & 0.5 & 2 & 35.8 & 35.8 & 35.7 & 36.0 & 37.0 & 36.8 \\
yeast & 6mer & 0.15 & 2 & 35.8 & 36.0 & 35.8 & 36.3 & 36.0 & 37.0 \\
yeast & 6mer & 0.5 & 2 & 36.2 & 35.5 & 35.3 & 35.9 & 36.5 & 36.1 \\
yeast & char & 0.15 & 2 & 36.3 & 36.4 & 35.1 & 36.1 & 36.8 & 36.9 \\
yeast & char & 0.5 & 2 & 36.2 & 35.3 & 35.3 & 36.0 & 36.7 & 36.9 \\
\end{longtable}
}

\section{Re-drawing the member/non-member canary split}

Which canaries are members and which are matched controls is fixed by one data seed, so within an
arm it is the same split for every model. With the mixed variant-derived canary set, the two groups of 50 records
differ by chance, and before any fine-tuning the released checkpoints already read the member group
as easier, by up to \NUM{0.9}\,\bits{} on a single draw. It is tempting to treat that as a fixed
offset and subtract it. Re-drawing the split four times shows that this is the wrong correction: the
offset is gone after a single epoch (\excess{} at round~1 is \NUM{0.00}$\pm$\NUM{0.05} on every
draw), a regression of each draw's round-50 \excess{} on its own round-0 offset has slope
\NUM{0.05}--\NUM{0.19}, and subtracting it in full inflates the across-draw spread. We therefore
retain the uncorrected statistic and use the round-0 reading as a diagnostic: it measures the
instability of a pretrained model's prior over two sets of variant-derived records, not a fixed property of the
canaries.

Table~\ref{tab:draws} gives the draws side by side, and the spread in it is not small: across the
three draws that admit a matched comparison, the same model's \excess{} moves by $0.06$--$0.27$\,
\bits{} at a single copy and by $0.07$--$0.64$ at sixteen, with the 5-mer the least stable. The
fourth draw admits none at all: its models' overfitting branches do not overlap, so there is no
common held-out loss to read them at. This is the reason the human arm of the main text is reported
with a never-trained control in every cell rather than with a single-draw error bar.

\begin{table}[htbp]
\centering\small
\caption{\emph{Generated.} The same matched-loss measurement under four member/non-member splits,
\excess{} \bits{} at the highest matched loss of each draw.}
\label{tab:draws}
% Generated by pilot2/appendix_tables.py -- do not edit. A11: the four canary-split draws.

\begin{tabular}{lcccccc}
\toprule
& \multicolumn{3}{c}{1 copy} & \multicolumn{3}{c}{16 copies} \\
\cmidrule(lr){2-4}\cmidrule(lr){5-7}
model & seed 1234 & seed 2345 & seed 3456 & seed 1234 & seed 2345 & seed 3456 \\
\midrule
DNABERT 3-mer [causal] & 0.44 & 0.36 & 0.37 & 1.55 & 1.41 & 1.37 \\
DNABERT 4-mer [causal] & 0.07 & 0.23 & 0.34 & 1.21 & 0.88 & 1.26 \\
DNABERT 5-mer [causal] & 0.06 & 0.14 & 0.17 & 1.21 & 0.57 & 0.93 \\
DNABERT 6-mer [causal] & 0.26 & 0.49 & 0.38 & 1.46 & 1.28 & 1.34 \\
DNABERT-2 BPE [causal] & 0.19 & 0.29 & 0.18 & 1.61 & 1.20 & 1.35 \\
hyena\_medium-160k & 0.34 & 0.30 & 0.28 & 1.52 & 1.47 & 1.45 \\
\midrule
\multicolumn{7}{l}{\footnotesize Draw 4567 has no matched comparison: its models' overfitting branches do not overlap.} \\
\bottomrule
\end{tabular}

\end{table}

% ---- The two sections moved out of the main text on 28 October for the ten-page
% limit. They are input from sections/, which holds the only copy of their text.
% MOVED OUT OF THE MAIN TEXT on 28 Oct for the ten-page limit (the author's decision;
% notes/22_page_cut_plan.md). This file is now input by appendix.tex, not by main.tex,
% which is why its heading is \section and its main-text references are numbers.
% Privacy-first section 4: recovery near each run's best saved held-out loss.
% Budget 1.75 pages including Table~\ref{tab:cost} and Figure~\ref{fig:cost}.
% Every number here is from the single-basis re-measurement of 19-23 Oct; see notes 3.58-3.59.
\section{Sequence recovery near the best saved held-out loss}
\label{sec:leak}

The question a practitioner faces is whether leakage appears at a checkpoint they would choose,
not whether a model memorizes after being trained past the point of usefulness. We
therefore pair every leakage rate with its held-out loss cost relative to the
best checkpoint that run itself produced. The decoding protocol and the number of record insertions jointly determine whether reproduction
occurs (Figure~\ref{fig:overview}b, upper half).
The second quantity makes the first operationally meaningful,
because it states how much predictive loss accompanies the measured recovery.

\paragraph{The cheapest leak costs seven thousandths of a bit.}
Table~\ref{tab:cost} orders the released checkpoints by that price. DNABERT-2 reproduces 30\% of
planted records at \textbf{+0.007\,\bits} above its own optimum. Nucleotide Transformer gives up
\textbf{+0.03\,\bits} to leak 8--14\% across three seeds. These checkpoints remain close to the best saved held-out loss, despite measurable recovery.

The never-trained control remains at zero throughout these evaluations. Records the model never saw are reproduced 0\% of the time in every
cell, every tier and every seed, and the same decoder scores the never-trained window at the foot of
each block on the same host, so the zero is a property of the measurement and not of a
separate baseline.

\begin{table}[tbp]
\centering
\caption{Released masked checkpoints on \emph{E.~coli}, records planted $16\times$, three seeds where shown. NT uses 50 records per tier;
the DNABERT-2 round-30 prefix and search cells use 10 and 20 records, respectively. \emph{Cost} is held-out \bits{} above that run's own best saved
checkpoint. \emph{Standard} is the prefix-conditioned protocol the field uses; \emph{ours} is the
re-masking decoder of \S\ref{sec:measure}. The DNABERT-2 rates are 3/10 and 0/20,
so that row compares protocols with different record coverage.}
\label{tab:cost}
\begin{tabular}{llccc}
\toprule
checkpoint & corpus & cost & standard & ours \\
\midrule
DNABERT-2 117M      & \emph{E.~coli} & \textbf{+0.007} & \textbf{30\%}      & 0\% \\
NT 250M (3 seeds)   & \emph{E.~coli} & \textbf{+0.03}  & 4 / 8 / 4\%        & \textbf{8 / 8 / 14\%} \\
NT 250M (3 seeds)   & \emph{E.~coli} & +0.44           & 26 / 26 / 30\%     & \textbf{66 / 68 / 76\%} \\
NT 500M             & \emph{E.~coli} & +0.46           & 36\%               & \textbf{78\%} \\
\midrule
never-trained control & --- & --- & 0\% & 0\% \\
\bottomrule
\end{tabular}
\end{table}

\paragraph{The two protocols disagree, and neither dominates.}
Comparisons across Table~\ref{tab:cost} show why this measurement cannot rely on a single extraction
test. On fixed-width non-overlapping tokenizers the standard protocol under-reports by
roughly a factor of two: at +0.44\,\bits{} it recovers 26--30\% where ours recovers 66--76\%. On
DNABERT-2 the order reverses completely. The standard protocol recovers 30\% and ours recovers
nothing at all. \S\ref{sec:measure} explains this protocol dependence. For each model, we interpret the reported exposure
using the stronger tested protocol. These measurements therefore provide lower bounds on what an attacker could obtain and not artifacts of one decoder's blind spot.

\paragraph{A single copy of a human-variant-derived record can be recovered.}
Uniform random canaries provide a controlled test but omit the sequence structure
of genomic inputs. We therefore extend the measurement to chromosome-22 reference
windows carrying observed rare alleles, testing whether recovery also occurs beyond
uniform random strings. These are the variant-derived records of \S\ref{sec:measure}. This is the original
mixed collection: 145 of 400 windows are chromosome-compatible at the audited sites,
20 are within-person genotype composites, and 235 are unobserved composites
(\S\ref{app:haplo}). Compatibility remains qualified at indel and multiallelic sites.
Recovery therefore demonstrates reproduction of a planted variant-derived record; it
does not establish full-base authenticity or individual linkage.

\begin{table}[tbp]
\centering
\caption{Nucleotide Transformer fine-tuned on human chromosome~22 with constructed rare-variant records
as canaries, including the mixed categories described above. \emph{1 copy} means one record appearing once in the
fine-tuning set. The second column is held-out \bits{} at that checkpoint. Round 50 is three seeds,
decoded in one batch; the earlier rounds are the first seed only, and are there to show the
trajectory rather than to carry an interval.}
\label{tab:human}
\begin{tabular}{lccccc}
\toprule
& & \multicolumn{2}{c}{16 copies} & \multicolumn{2}{c}{1 copy} \\
\cmidrule(lr){3-4}\cmidrule(lr){5-6}
checkpoint & held-out \bits & standard & ours & standard & ours \\
\midrule
NT 250M, round 15          & 1.411         & 6\%      & 6\%                & 0\%    & 0\% \\
NT 250M, round 30          & 1.501         & 10\%     & 28\%               & 0\%    & 0\% \\
NT 250M, round 50, 3 seeds & 1.632--1.648  & 16--18\% & \textbf{48--54\%}  & 0--2\% & \textbf{8\%} \\
NT 500M, round 15          & 1.372         & 16\%     & 14\%               & 0\%    & 0\% \\
NT 500M, round 50, 3 seeds & 1.668--1.788  & 24--50\% & \textbf{58--72\%}  & 0\%    & \textbf{12--18\%} \\
\midrule
never-trained control & --- & 0\% & 0\% & 0\% & 0\% \\
\bottomrule
\end{tabular}
\end{table}

Table~\ref{tab:human} presents the central result of this section. A variant-derived
record, present exactly once in the fine-tuning data, is recovered exactly by the 250M model in 8\% of cases in \emph{every one} of three seeds, and by the 500M in
12--18\%, between one case in twelve and one in six, after our fine-tuning of publicly downloadable starting weights, with the
matched never-trained control at zero in all of them. The repeated tier
behaves as the \emph{E.~coli} arm does, so the single-copy column is not an artifact of a canary
appearing often enough to be learned as vocabulary.

\paragraph{Conditional prediction and measured recovery.}
For DNABERT-2 on chromosome~22, single-copy records at the final checkpoint have
93--95\% nucleotide-level LOO accuracy and a power summary of 0.21--0.39 across
three seeds, while both tested decoders recover none of their hidden targets
(Table~\ref{tab:loo}). LOO supplies every other true target nucleotide; these are
conditional-prediction diagnostics, not bounds or attainable whole-record recovery
rates (\S\ref{sec:measure}).
Separately, for the \emph{E.~coli} sixteen-copy condition in Table~\ref{tab:cost},
DNABERT-2's prefix protocol recovers 30\% whereas the tested re-masking search
recovers 0\%. We report the stronger measured protocol for that condition.
Neither search failure alone establishes an absence of record-specific learning.

Table~\ref{tab:loo} compares the listed released-model checkpoints on chromosome~22, because the
same caution distinguishes conditional prediction with true target context from whole-record
decoding without that context. Read down the single-copy rows: the overlapping DNABERT family is the
case with the smallest power summary, $0.01$--$0.02$, alongside zero tested recovery.
Elsewhere the power summary is 0.14 to 0.41 while measured recovery is 0--18\%.
These quantities describe different tasks: the probe receives the other true target
positions, whereas the decoder must infer them. No power value here demonstrates that
a better decoder can attain it or excludes recovery by another protocol.

\begin{table}[tbp]
\centering
\caption{Conditional prediction and tested recovery on chromosome~22, 50 records per tier,
three seeds where shown. LOO supplies all other true target positions. \emph{Power diag.}
is the mean of per-record $p_i^{s_i}$, not a recovery bound or the power of mean accuracy.
NT probes tokens; DNABERT-2 and overlapping DNABERT probe 48 nucleotides. Control powers
are $0.000$--$0.003$. We rescored the overlapping-family prefix column at all fifty
records (\S\ref{app:prov}). Protocols and budgets are specified in \S\ref{sec:decoders}.}
\label{tab:loo}
\begin{tabular}{llcccc}
\toprule
checkpoint & copies & LOO top-1 & power diag. & standard & ours \\
\midrule
NT 250M, round 50 (3 seeds) & 16 & 0.96--0.97 & 0.81--0.85 & 16--18\% & \textbf{48--54\%} \\
                       & 1  & 0.57--0.66 & 0.14--0.23 & 0--2\%   & \textbf{8\%}  \\
NT 500M, round 50 (3 seeds) & 16 & 0.96--0.99 & 0.84--0.92 & 24--50\% & \textbf{58--72\%} \\
                       & 1  & 0.67--0.80 & 0.24--0.41 & 0\%      & \textbf{12--18\%} \\
DNABERT-2, final (3 seeds) & 16 & 0.99--1.00 & 0.85--0.91 & \textbf{40--44\%} & 0\% \\
                       & 1  & 0.93--0.95  & 0.21--0.39  & 0\%     & 0\% \\
DNABERT $k=3$--6, round 120 & 16 & 0.99--1.00 & 0.70--0.87 & 0--50\%  & 0--4\% \\
                       & 1  & 0.61--0.66  & 0.01--0.02  & 0\%     & 0\% \\
\midrule
never-trained control  & 0  & 0.07--0.51  & 0.000--0.003 & 0\%     & 0\% \\
\bottomrule
\end{tabular}
\end{table}


% MOVED OUT OF THE MAIN TEXT on 28 Oct for the ten-page limit (the author's decision;
% notes/22_page_cut_plan.md). This file is now input by appendix.tex, not by main.tex,
% which is why its heading is \section and its main-text references are numbers.
% Privacy-first section 7: what a practitioner can do, and why the usual metric will not help.
% Budget 1.25 pages including Table~\ref{tab:blind} and Figure~\ref{fig:blind}.
% Early stopping: outputs/nt_{ckpt,seeds}/nt250m_mlm_ecoli_s*, ep3/ep15/ep50, one basis 19 Oct.
% Blind metric: outputs/nt_ckpt/nt250m_mlm_human_s0/utility_at_leak.json and bpe_human/*, 19-20 Oct.
\section{What a practitioner can do, and what will not warn them}
\label{sec:blind}

Our measurements inform checkpoint selection in two ways: early checkpoints can have
lower measured recovery, and downstream task scores need not track that change. The comparisons below
are specific to the model, checkpoint and copy tier; they do not establish early stopping as a general defense.

\paragraph{Early checkpoints can resemble controls on the tested diagnostics.}
At round~3 of NT~250M on \emph{E.~coli}, seeds~1 and~2 have zero prefix and re-masking
recovery at both one and sixteen copies (50 records per tier). Single-copy target surprise
is $2.168$ and $2.154$\,\bits{}, against $2.112$ and $2.110$ for the respective controls;
LOO top-1 accuracy is $0.00$ after rounding, against $0.01$ for controls. At sixteen copies,
target surprise is $2.128$ and $2.115$\,\bits{} and LOO top-1 is $0.01$ in both seeds.
These are token-level conditional probes, not tests that a record is absent. The near-control
diagnostics and zero recovery describe these checkpoints under the tested protocols.

For records planted sixteen times, a further increase of three hundredths of a bit
in held-out loss yields 8\%, 8\% and 14\% exact recovery across three seeds. At a larger increase of +0.44\,\bits, the recovery rate reaches 66--76\%. Earlier stopping reduces recovery in this arm, with a sharp change between saved checkpoints. Our own from-scratch backbone tolerates roughly ten times that drift before
showing exact recovery under the tested protocol, so a study of that backbone alone reads more reassuringly than
released weights warrant. That is why this paper is built on released weights.

\paragraph{The downstream selection metric changes in both directions.}
Checkpoint selection can instead use downstream task scores.
Table~\ref{tab:blind} puts that score beside the leak at selected checkpoints of the same runs.

\begin{table}[tbp]
\centering
\caption{Downstream task score and leak measured on the same weights at selected checkpoints.
GUE promoter classification, linear probe on frozen mean-pooled states, Matthews correlation on the
full held-out split. \emph{id} is single-copy identification in the record's own host window, as in
Table~\ref{tab:ident}; an earlier version of this table quoted the average over the two host
conditions, which is a different quantity. Identification is over three seeds. DNABERT-2 MCC
entries give the three-seed minimum--maximum; NT MCC is from seed~0.
The DNABERT-2 best-checkpoint label uses mean linear-probe MCC across seeds.}
\label{tab:blind}
\begin{tabular}{llcccc}
\toprule
model & checkpoint & held-out \bits & MCC & 16-copy verbatim & 1-copy id \\
\midrule
NT 250M, chr22 & round 15 (best MCC) & 1.411 & \textbf{0.802} & \textbf{6\%}  & \textbf{0.04--0.08} \\
               & round 30            & 1.501 & 0.798          & 10\%          & 0.30--0.34 \\
               & round 50            & 1.632 & 0.780          & 18\%          & 0.60--0.66 \\
\midrule
DNABERT-2, chr22 & \shortstack[l]{round 15\\(best mean MCC)} & 1.264 & 0.798--0.801 & 0--2\%   & 0.04 \\
                 & round 120           & 1.600 & 0.784--0.792 & 42--46\% & 0.74--0.80 \\
                 & final               & 1.729 & 0.782--0.787 & 40--44\% & 0.92--0.96 \\
\bottomrule
\end{tabular}
\end{table}

On DNABERT-2 the score declines across the run by 1.1--1.6 points of MCC in total, monotonically in
two seeds and with one reversal at round~50 in the third, while single-copy identification climbs
from 4\% at round~15 to 92--96\% at the final checkpoint in its own window
(Figure~\ref{fig:blind}); matched own-window controls at round~15 are 0--2\%.
Selecting by mean linear-probe MCC across the three seeds chooses round~15, where measured
own-window membership rates are lower than at later checkpoints. Individual-seed selection
chooses round~15 for seeds~0--1 and round~50 for seed~2. This linear-probe score is not
a privacy signal: across the entire interval in which the model learns to disclose
membership, a practitioner watching it sees about a point of MCC of drift, which is what an
ordinary run looks like.

On NT fine-tuned on chromosome~22 the situation is worse, because the metric does not select a checkpoint with zero measured leakage. The highest task score belongs to the earliest checkpoint saved, and that one
already reproduces 6\% of repeated records in seed~0; single-copy own-window
rank-1 rates are 4--8\% across the three seeds. Choosing the best
available model by downstream performance means choosing a model that is already leaking.

\begin{figure}[!b]
\centering
\includegraphics[width=\linewidth]{F7_blind.pdf}
\caption{Downstream score and leak, measured on the same weights at the same checkpoints.
\textbf{Left:} DNABERT-2 on chromosome~22: the score falls by 1.1--1.6 points of MCC while
single-copy own-window rank-1 rates rise from 4\% at round~15 to 92--96\% at the final
checkpoint; round-15 matched controls are 0--2\%. Solid and dashed curves show the three-seed
means of MCC and rank-1, respectively; shading spans minimum--maximum MCC. The marked best
checkpoint maximizes mean linear-probe MCC.
\textbf{Right:} Nucleotide Transformer 250M on the same corpus, seed~0. The best downstream
score is at the earliest checkpoint shown, where single-copy own-window rank-1 is 8\%. Linear
probe on frozen mean-pooled states, GUE promoter classification, full held-out split.}
\label{fig:blind}
\end{figure}



\paragraph{A larger background corpus dilutes the leak, and does not prevent it.}
The obvious question about every number above is whether a real cohort, which is larger than ours,
is safer by being larger. Our experiment varies the background corpus rather than the participant cohort: the same backbone, cut
size, objective, canaries and schedule, with the fine-tuning corpus at 1{,}000, 2{,}500, 5{,}000,
10{,}000 and 20{,}000 windows of chromosome~22, two seeds each. The number of participants, how
many windows each contributes and how closely they are related are fixed by construction and were
not manipulated, so this experiment measures window-count scaling and bears on cohort size only to the
extent that the two move together. Table~\ref{tab:scale} reports results at matched held-out loss. This is necessary because, at a
fixed epoch count, a larger corpus also entails more optimization updates.

Dilution is real and it is not a defense. Twenty times the corpus divides single-copy verbatim
recovery by about four at $2.05$\,\bits{} (97\% to 23\%) and takes identification from $1.00$ to
$0.64$. At $2.20$\,\bits{}, all three corpus sizes that reach that loss have rank-1 rates
of $0.91$ or better; at $2.40$, every size that reaches it is at $1.00$. A larger corpus shifts the onset of the same leakage along the loss axis. This delay is useful and
uses the same measure as \S\ref{sec:leak}, but it does not provide immunity. At each corpus size's own held-out optimum, meanwhile,
single-copy verbatim recovery is 0\% at every size and sixteen-copy recovery is 78--100\% at every
size: the stopping rule does the work, not the size of the corpus.

\begin{table}[tbp]
\centering
\caption{Identification and verbatim recovery for one planted record as the background corpus grows,
read at matched held-out loss, mean of two seeds, records planted once. The manipulated quantity is the number of
chromosome~22 windows, not the number of participants. A dash means that corpus never reaches that
loss in 30 epochs. \emph{Best loss} is the corpus's own optimum, which improves with size as it
should.}
\label{tab:scale}
\begin{tabular}{lccccc}
\toprule
windows & best loss & id @2.05 & id @2.20 & verbatim @2.05 & verbatim @2.20 \\
\midrule
1{,}000  & 1.954 & 1.00          & ---           & 97\%          & --- \\
2{,}500  & 1.910 & 1.00          & ---           & 96\%          & --- \\
5{,}000  & 1.883 & 0.95          & 1.00          & 56\%          & 94\% \\
10{,}000 & 1.789 & 0.74          & 1.00          & 50\%          & 82\% \\
20{,}000 & 1.708 & \textbf{0.64} & \textbf{0.91} & \textbf{23\%} & \textbf{37\%} \\
\bottomrule
\end{tabular}
\end{table}

\paragraph{Recommendation.}
These findings inform four practical recommendations, presented as guidance distinct from the
experimental results themselves.

\emph{Stop where held-out loss stops improving, and compute it.} In this trajectory, held-out loss worsens before the observed rise in exact
recovery, and it is inexpensive to monitor during fine-tuning. This stopping point does not define a sharp boundary between privacy and leakage. At that checkpoint single-copy identification is $0.04$ and exact recovery is zero,
but the decoder's output is already closer to a planted single-copy record than to a never-trained
one in the seed-averaged descriptive comparison ($31.7$ against $33.2$ nucleotides
wrong out of 48); the within-cell gaps are small relative to their standard errors
(\S\ref{sec:ident}). The earlier checkpoint has lower measured leakage under the tested protocols;
this does not establish that its records are unreachable by another attack. It is a measured
risk reduction in these runs, not a safety property of a release: at that same checkpoint a
record planted sixteen times is already ranked first.

\emph{Do not treat downstream score as a safety signal.} Between each model's own held-out optimum
and the end of its run, an interval over which extraction goes from 0--6\% to 36.7--100\%, the available frozen-probe means
on six tasks across two organisms show maximum absolute changes of $0.024$ MCC on the four DNABERT
arms and $0.083$ on the other two, while the fine-tuned head moves \emph{upward}
on four of the six tasks for DNABERT-2, telling a practitioner who watches it to keep training.
Separately, the frozen promoter probe on chromosome~22 declines by 1.1--1.6 MCC points
as membership rank-1 rates rise. Across these evaluations, downstream utility does not provide
a consistent warning of the observed leakage.

\emph{Evaluate masking under the intended model and release conditions.} On our own
backbone, trained on the same corpus for the same budget, replacing next-token training with masked
training takes verbatim extraction from 100\% to \emph{zero} in every cell we ran (four corpora,
two mask rates, three cut sizes), and it \emph{improves} held-out loss, from 2.3--4.2\,\bits{} to
1.66--1.98. This is a 100-percentage-point reduction in exact extraction within the
controlled objective comparison, for the tested backbone and decoding protocols.
The masked pretrained models in Table~\ref{tab:cost} also exhibit leakage, which limits
transfer of this observation. That comparison changes architecture, scale and training
history together; it does not isolate pretraining as the cause or establish masking as
a privacy defense. Both leakage endpoints require evaluation on the intended released model.

\emph{Evaluate tokenization within the intended backbone and objective.} Under the drawn
placement used in the controlled sweep, a model whose canary occupies 24 or 48 content slots
returned no record whole at any loss we drove it to, including losses far past usefulness. That
sweep measured exact extraction as its endpoint; whether more slots also reduce membership
disclosure was not tested there. The placement condition is also essential to this interpretation: at a \emph{fixed} offset the character tokenizer reaches 100\% at $2.20$\,\bits{}
(\S\ref{sec:basis}, Table~\ref{tab:capacity}). The effect therefore depends on placement,
and the differing target surprise does not support equal-learning comparisons. HyenaDNA
and the failed overlapping-family prediction further limit a general slot-count explanation
(\S\ref{sec:slots}). The magnitude remains substantial: exact recovery rises from 0\% with 2-mers
to 92\% with 6-mers at matched loss $2.40$\,\bits{} in the sixteen-copy sweep. Fine-grained
tokenizers carry their own costs in sequence length and downstream accuracy, which depend on the task and were not evaluated comprehensively here. Slot count is therefore a
design choice whose costs merit measurement; we do not recommend it without qualification.


\section{Measurement provenance and corrections}
\label{app:prov}
The following corrections determine which measurements the paper reports; detailed
development history is retained in \path{notes/17_experiment_log_week3_EN.md}.
\begin{itemize}\itemsep2pt
\item \textbf{Record count:} the DNABERT-2 trajectory and overlapping-family prefix
trend were rescored at 50 records per cell rather than the original ten. The former
had appeared two checkpoints late; the latter initially read $0/0/10/40\%$.
Tables~\ref{tab:ident} and \ref{tab:loo} and \S\ref{sec:slots} use the full-count results.
\item \textbf{Controls and replication:} the 18 DNABERT-2 cells reuse 50 records across
six checkpoints and three seeds, so no pooled independent-trial rate is reported.
We retrained the HyenaDNA and overlapping-DNABERT controls because the original
weights were unavailable; \S\ref{app:bh} checks the reproduced member/loss values.
\item \textbf{Protocol and access:} NT~500M's single-copy zero applies to prefix
decoding, while re-masking yields $0.12$--$0.18$. Table~\ref{tab:blind} uses own-window
membership rather than an average of own and fresh conditions.
\item \textbf{Diagnostic interpretation:} legacy \texttt{loo\_bound} values are retained,
but no longer claimed to certify recovery bounds (\S\ref{app:measurement}). The
BPE normalization audit and corrected-scorer coverage remain separate
(\S\ref{app:bpenorm}).
\end{itemize}

\section{Experimental configurations and table interpretation}
\label{app:config}

Table~\ref{tab:config} reports each arm's settings directly from the run's recorded
\texttt{args.json}, rather than from the manuscript text.

\begin{table}[htbp]
\centering\small
\caption{One row per arm. The masked arms differ in scheme as well as rate: DNABERT-2 and
Nucleotide Transformer mask tokens independently, while the overlapping family masks spans of $k$
consecutive tokens with starts drawn at rate$/k$, at 0.15, 0.20, 0.25 and 0.30 for $k=3,4,5,6$.
\emph{Checkpoints} lists the union of principal reported snapshots across runs; supplementary
experiments retain their separately stated configurations. The released checkpoints are fine-tuned,
not trained from scratch; the last row is the only backbone we train.}
\label{tab:config}
\begin{tabular}{llllll}
\toprule
arm & tokenizer & objective & lr & epochs & checkpoints \\
\midrule
DNABERT-2, 117M      & byte-pair          & masked, 0.15 & \num{2e-5} & 150 & 15, 30, 50, 80, 120, 150 \\
DNABERT, masked arm & overlapping \kmer{k} & span, 0.15--0.30 & \num{1e-4} & 150 & 15, 30, 50, 80, 120, 150 \\
DNABERT, causal arm & overlapping \kmer{k} & causal & \num{2e-5} & 50 & 1, 3, 8, 15, 30, 50 \\
NT 250M and 500M     & fixed-width 6-mer  & masked, 0.15 & \num{2e-5} & 50  & 3, 15, 30, 50 \\
HyenaDNA medium-160k & single nucleotide  & causal       & \num{2e-5} & 50  & 1, 3, 8, 15, 30, 50 \\
controlled backbone  & $k=1$--6, spaced   & next-token   & \num{3e-4} & 30  & 1, 2, 3, 5, 8, 12, 20, 30 \\
\bottomrule
\end{tabular}
\end{table}

\paragraph{Principal campaign settings and supplementary exceptions.}
On chromosome~22, NT~250M saves rounds 15, 30 and 50, while NT~500M saves rounds 3, 15 and 50.
On \emph{E.~coli}, NT~250M seed~0 saves rounds 15, 30 and 50, seeds~1 and~2 save rounds 3, 15 and 50, and NT~500M seed~0 saves rounds 3, 15 and 50.
The principal released-model campaigns use five thousand host windows, canaries planted
1, 4 and 16 times with 50 records per tier, each copy
at an independently drawn offset; 50 never-trained control records; a candidate pool of 100 real
segments; \texttt{data\_seed} 1234, so all seeds of an arm share one member, control and pool
split and re-score the same records; batch 8 with gradient accumulation 2. The from-scratch sweep is an exception, using larger evaluation sets and batches where
computationally affordable: 100 records per tier, 100 controls, a pool
of 500, batch 64, on uniform rather than human sequence. Supplementary controls use
their stated settings: Table~\ref{tab:capacity} has fixed placement and 20 records
per tier, and the edge-conditioned audit (Table~\ref{tab:a10}) also uses 20 records
per tier. The large-corpus experiment above uses 100{,}000 windows.

\paragraph{Decoders and their budgets.}
\S\ref{sec:decoders} lists the update units, conditioning and restart/pass budgets.
We report both applicable protocols, identifying the stronger measured result; LOO
diagnostics are reported separately because they use additional true target context.

\paragraph{Reading at matched loss.}
Within each seed, rates are linearly interpolated between neighboring checkpoints in held-out-loss
order that bracket the target; available seed estimates are then averaged. The controlled sweeps
use numbered epochs 1, 2, 3, 5, 8, 12, 20 and 30. The separately labeled \texttt{final}
evaluation duplicates epoch 30 and is counted only once. A cell is left blank where the target
lies outside that run's range; there is no extrapolation. Where a leak is priced
in \bits{}, the baseline is that run's own best saved checkpoint, so no price is comparable across
arms except through this rule.

\paragraph{A release check a group can run, and what it does not certify.}
The inputs are ones the releasing group already holds and an outside adversary does not: the
fine-tuning corpus, the manifest of which windows went into it, and the weights at each saved
checkpoint. From those it assembles a candidate pool of real segments drawn from the same
distribution and absent from the corpus, and a set of never-trained control records placed in the
members' host windows and offsets for own-window controls; fresh controls use independent
background and offset draws. Then, at each checkpoint, it reads two
numbers in each of two access conditions (the record's own training window, and a window the
model never saw) and the exact-recovery rate of the decoder matched to that model's vocabulary.
What it reports is the member rate beside the control rate on the same cells, per condition, per
decoder, with the record count beside each. Our code runs exactly that
(\S\ref{sec:commands}).

Three limitations accompany these measurements. The own-window numbers are a privileged-context diagnostic:
an adversary holding the corpus could search the corpus and would not need the weights, so only
the fresh-window pair speaks to what a released checkpoint gives up. The statistic is a ranking
against a candidate set: rank~1 defines a decision rule, but does not supply a
transferable low-false-positive-rate threshold or a calibrated per-record error rate
(\S\ref{sec:ident}). And a checkpoint that reads clean has been shown clean against these
decoders, these budgets and this pool, which is a statement about the tested modes and not a
safety certificate.

\section{Exact commands}
\label{sec:commands}

Every command runs from the repository root. The commands reproduce the reported
measurements, figures and appendix tables. Each restart count below is the value
recorded in the result file that command writes; the earlier \texttt{recon\_bpe\_*} and
\texttt{recon\_*r32*} runs these supersede used 32 restarts.

\begin{footnotesize}
\begin{verbatim}
# One cell of the 23 Oct batch, per vocabulary (Tables A12, A13 and A6 of this
# appendix). The batch is these three commands over 19 (run, checkpoint) pairs,
# one GPU each, logged to outputs/chr22_rerun. DB5 is the 5-mer run directory.
NT=outputs/nt_ckpt/nt250m_mlm_human_s0
DB5=outputs/dnabert_long/dnabert5_nt0.05_lr0.0001_human_s0
BPE=outputs/bpe_human/dnabert2bpe_human_s0

python -m pilot2.recon_audit   --root $NT  --ckpt ckpt_ep50.pt \
    --restarts 128 --max_probes 50 --name rb_ep50
python -m pilot2.recon_overlap --root $DB5 --ckpt ep120.pt \
    --restarts 16  --max_probes 50 --name rb_ep120
python -m pilot2.recon_bpe     --root $BPE --ckpt ep120.pt \
    --restarts 16 --max_probes 50 --score pernt --name rb_ep120

# the tokenizer sweep, k = 1..6, three seeds, window 300 (per-seed Table A4)
# sweep.sh activates the recorded conda environment glm.
EXTRA="--epochs 30 --tiers 1,4,16 --probe_offset random"
EXTRA="$EXTRA --probes_per_tier 100 --n_nonmember 100"
EXTRA="$EXTRA --save_epochs 1,2,3,5,8,12,20,30 --score_snapshots"
PILOT2_WINDOW=300 bash pilot2/sweep.sh outputs/kcurve300 "ar" "0 1 2" \
    "$EXTRA" "char 2mer 3mer 4mer 5mer 6mer"
python -m pilot2.kcurve_report --root outputs/kcurve300


# re-score a saved checkpoint and compare with its original scores
python -m pilot2.rescore --root outputs/bpe_human/dnabert2bpe_human_s0 \
    --ckpt ep120.pt --name scores_ep120_rb --compare scores_ep120.json

# the published derived table, the figures, and this appendix's tables
python -m pilot2.export_figure_data --root outputs --out results/figure_data.csv
python -m pilot2.figures_privacy --out paper/figures
python -m pilot2.figures_main10 --out paper/figures
python -m pilot2.appendix_tables --out paper/tables
\end{verbatim}
\end{footnotesize}

\noindent
\texttt{GLMEXT\_DATA} must be exported in the shell that runs these: it is set in \texttt{.bashrc},
which a non-interactive SSH session does not read, and a batch that inherits it from an interactive
shell will fail instantly in every cell when it is re-run from a script.


\section{Measurement definitions and decoder diagnostics}
\label{app:measurement}
This section expands the definitions in \S\ref{sec:measure}; the score, rank,
re-masking update and common-interval construction remain in the main text.
The historical BPE
search files for the main chromosome-22 remeasurement are \texttt{rb\_ep30.json},
\texttt{rb\_ep120.json} and \texttt{rb\_epfinal.json} for each of three seeds,
with 16 restarts and 50 records per evaluated tier. The separate corrected scorer
audit is seed 0's \texttt{rbfix\_epfinal.json}: 60 cells in total, with 16 restarts.
It documents zero exact recovery after correcting the divisor in those audited cells,
not a corrected-scorer sweep over all seeds, contexts or checkpoints.
\paragraph{The held-out loss scale.}
We express the cost of each leak in \bits{} above the lowest held-out loss among that run's saved
checkpoints. For token-based scorers, held-out loss is the teacher-forced surprise of a held-out window divided by its
nucleotide count. This normalization removes the reduction in reported loss that coarser tokenization
would obtain simply from fewer predicted slots:
%
\begin{equation}
\label{eq:heldout}
\ell(M)=\frac{1}{N}\sum_{i=1}^{N}\frac{-1}{W}
\sum_{j\in\mathcal{J}(h^{(i)})}\log_{2}P_{M}\bigl(u_{j}\mid c_{j}\bigr),
\qquad
\ell^{\star}=\min_{M'\in\mathcal{K}}\ell(M'),
\end{equation}
%
over the first $N$ windows of a held-out split with no inserted canaries, where $\mathcal{K}$ is the
set of snapshots the run saved and $c_{j}$ is the scoring context: the prefix
for a next-token model, and every other slot for DNABERT-2 and NT.
The DNABERT-2 chromosome-22 arm evaluates all 200 test windows;
the controlled synthetic sweeps evaluate all 500 test windows. Overlapping DNABERT
uses a different estimator: \texttt{heldout\_floor} in \texttt{pilot2/dnabert.py} averages
nucleotide surprise at positions $k,k+7,\ldots<W-k$, then averages windows.
Its ranking score uses positions $o,o+3,\ldots<o+96$. The 50-epoch
\texttt{human\_real3} causal runs supply the cross-family membership rates; their
\texttt{human\_bh} control retraining uses the same objective.
\texttt{nt\_logprobs\_causal} applies causal attention and predicts the token starting
at $j-k+1$, restricting its preceding $k-1$ letters and normalizing over four continuations
(positions $j<k-1$ are excluded). The separate 150-epoch \texttt{dnabert\_long}
masked runs supply the decoder/LOO comparison. There \texttt{nt\_logprobs} masks
every covering token and reads token $\min(j,W-k)$; consistent $k$-mer logits are
marginalized onto the letter and normalized over A/C/G/T. These nucleotide scores
replace the token sum for this family; the two campaigns do not share weights.
Weighting each window equally
rather than pooling slots is what makes $\ell$ an axis on which a \kmer{6} cell and a character
cell can be compared at all, and $\ell(M)-\ell^{\star}$ is the price every later table quotes.

\paragraph{The matched own-window control.}
A never-trained record has no training row. We therefore assign each control a member's row and
offset, and the ranking pass overwrites $[o,o+2L)$ with \emph{every} candidate in turn. The memorized string is therefore not itself in $C$, and the two conditions
%
\begin{equation}
\label{eq:control}
\bigl(x\in\mathcal{R}_{\rho>0},\ h_{\mathrm{own}},\,o\bigr)
\quad\text{against}\quad
\bigl(x'\in\mathcal{R}_{0},\ h_{\mathrm{own}},\,o\bigr)
\end{equation}
%
differ in membership alone. A control near chance argues against a generic training-window advantage as the sole
explanation for the member rates; it does not separate a model whose
conditional at that offset is sharp on a string absent from $C$, which would drive the control to
chance as well.

\paragraph{Conditional-prediction diagnostics.}
The leave-one-out (LOO) probe supplies every other target slot at its true value and
takes the argmax at one. For record $i$, let $p_i$ be its per-slot top-1 rate and $s_i$
the number of probed slots. We report both the mean accuracy and a per-record power summary:
\begin{equation}
\label{eq:ceiling}
\widehat b=\frac{1}{n}\sum_{i=1}^{n}p_i^{\,s_i},
\qquad \widehat p=\frac{1}{n}\sum_{i=1}^{n}p_i,
\qquad s_i\equiv s\ \Longrightarrow\ \widehat b\ge\widehat p^{\,s}.
\end{equation}
The inequality is Jensen's inequality. Neither summary is a recovery probability:
$p_i$ measures prediction given the other true target positions, not the probability of
a correct step during free decoding. The probe is token-level for NT and the non-overlapping
masked backbone, and nucleotide-level ($s_i=48$) for DNABERT-2 and overlapping DNABERT.

A limited abstract implication holds for \emph{complete-state reconstruction}: if a decoder
returns a fixed point of exactly the probe's update map, with identical constraints,
tie-breaking and numerical evaluation, then recovery of the full true state requires
$p_i=1$, and its indicator is at most $\mathbf1[p_i=1]\le p_i^{s_i}$. This implication
does not certify the reported experiments. In \texttt{recon\_audit.py}, covering tokens
can extend beyond the hidden 48\nt{}; leading revealed letters are constrained but
trailing host letters within the final token are not. Exact recovery scores only the
48\nt{} crop, whereas LOO checks the original full tokens. This mismatch affects NT
and masked-backbone runs using that routine. In addition, search and LOO use batched
and single-row forward passes, respectively, which need not produce identical argmaxes.

Saved outputs include exact recovery after early stopping with $p_i<1$: the NT~500M
human seed-0 round-50 single-copy probe~33 has $p_i=8/9$ and $p_i^9=0.346$.
Aligned counterexamples also occur, so boundary mismatch cannot explain every case;
the numerical cause has not been isolated by replay. Early stopping under the search
computation therefore does not certify the probe's necessary condition. The $95.8\%$
sub-cap statistic below describes selected-output stopping, not a verified bound.
No aggregate exceedance was observed in the six human seed-0 NT files checked, but
that ordering does not validate the failed recordwise premise. We use $\widehat b$
as a conditional-prediction diagnostic for \emph{all} reported arms. Historical JSON
keys named \texttt{loo\_bound} retain their values for reproducibility.

For single-copy DNABERT-2, $\widehat p=0.945$ and $\widehat b=0.268$, which exceeds
$\widehat p^{48}=0.065$ by a factor of $4.1$.
The BPE values refer to the historical scorer, with the common-interval correction
audited separately. Its zero exact recoveries do not validate a recovery bound.
The $93$--$95\%$ LOO accuracy establishes conditional predictive ability with true
target context; neither it nor the power summary establishes whole-record recoverability.

\paragraph{Teacher-forced covering-token surprise.}
Teacher-forced surprise sums log probabilities of the true covering tokens,
with revealed leading letters constrained, and divides by 48 target nucleotides.
Trailing host letters within the final token remain in the scored event. This is
a covering-token diagnostic normalized by target length, not the marginalized
probability of the hidden 48\nt{} alone.
Writing $c_j^{\mathrm{TF}}$ for the prefix with preceding target tokens supplied
at their true values, the diagnostic is
$T_{48}=-\frac{1}{48}\sum_{j\in\mathcal S}\log_2Q_{M,j}(u_j\mid c_j^{\mathrm{TF}})$,
with the content and revealed-prefix constraints used by the decoder.
This differs from the bidirectional LOO context and from a strict 48\nt{} event.

\paragraph{Slots.}
A decoder predicts complete slots. We therefore determine how many slots the $L$ hidden nucleotides
occupy, accounting for the record's offset within its host. Widening the target
to slot boundaries gives
%
\begin{equation}
\label{eq:slots}
s=\Bigl\lceil\frac{L+\varphi}{k}\Bigr\rceil
\;\overset{k\,\mid\,L}{=}\;\frac{L}{k}+\mathbf{1}\bigl[k\nmid o\bigr],
\qquad \varphi=(o+L)\bmod k,
\end{equation}
%
for a non-overlapping \kmer{k}; $s=L$ for the overlapping family and for single nucleotides, since
there a slot is a nucleotide. The diagnostic power $p^s$ in \eqref{eq:ceiling} decreases with
$s$ at fixed $0<p<1$. This mathematical comparison is not a recovery law and does not
establish equal target predictability across tokenizers. For a straddling token, exactness
requires only its hidden letters, whereas the token-level probe tests its full identity.
That is the account \S\ref{sec:slots} tests: whether recoverability tracks $s$ alongside a separate measurement of target surprise. A controlled sweep can show that it does; it cannot by itself exclude every other
way the tokenizer changes the fit, and we do not claim it has. The $L/k$ column of Table~\ref{tab:slots} is a nominal ratio, not an integer slot count.
Equation~\eqref{eq:slots} attains that ratio when $k$ divides both $L$ and $o$;
under the drawn placement those runs use,
the mean slot count is $48.0$, $24.5$, $16.7$, $12.8$, $10.4$ and $8.8$ for $k=1\ldots6$.

\paragraph{The decoders.}
The decoder implements the extraction attack. Consequently, the measured exposure of a checkpoint
depends on the decoder used for its vocabulary. The DNABERT-2 prefix protocol rolls the nominal 48\nt{} prefix back to the
preceding BPE boundary (41--48 revealed nucleotides in the chromosome-22 audit)
and retains the token-slot layout of the true window before masking identities.
We crop its saved predictions by nucleotide coordinates to the hidden 48\nt{}
for exactness and Hamming distance, counting missing letters as errors. This
uses a boundary-aligned input and a fixed 48\nt{} evaluation interval.
The \emph{standard} protocol generates one greedy continuation left to right from the
revealed half, taking the admissible-token argmax and feeding back its predictions, with no
sampling, beam search, refinement or restart. This entails one forward evaluation per generated
token (per nucleotide for overlapping DNABERT), not one forward pass per record. On a masked model it masks everything right
of the revealed prefix, the host suffix beyond the target included, and fills the target slots
sequentially while that suffix stays masked. This differs from the independent per-token masking
used to train DNABERT-2 and Nucleotide Transformer and the $k$-token span masking of the overlapping
family. Our \emph{re-masking} decoder instead treats the target's slots
$\mathcal{S}$ as latent. Writing $z=(z_{j})_{j\in\mathcal{S}}$ for the tokens placed on them,
$w(z)$ for the window carrying $z$ with everything outside $\mathcal{S}$ at its true value, and
$\mathcal{A}_{j}$ for the valid content tokens consistent with the revealed prefix letters straddled by slot $j$,
one sweep re-masks a single slot and takes the restricted argmax with every other slot held:
%
\noindent See \eqref{eq:remask} in the main text.

%
The sweep uses Gauss--Seidel rather than Jacobi updates: each slot conditions on its predecessors'
updated values. The iteration is therefore a deterministic map on $\prod_{j}\mathcal{A}_{j}$,
whereas parallel updates can oscillate. It runs for at most six sweeps from each of $R=128$
restarts, stopping when a sweep makes no changes. The attack must then pick a final state
\emph{without} the answer, which it does by the pseudo-log-likelihood of the state under the same
single-slot conditional used for its updates \citep{mireshghallah2022mlm}.
Here admissible content tokens are explicitly renormalized:
$Q_{M,j}(v\mid c)=\mathbf{1}[v\in\mathcal A_j]P_M(v\mid c)/
\sum_{a\in\mathcal A_j}P_M(a\mid c)$.
The denominator does not change a within-slot argmax but can change the ranking
of restart states; raw $P_M$ remains the probability used for membership scoring:
%
\begin{equation}
\label{eq:select}
\Lambda(z)=\sum_{j\in\mathcal{S}}\log_{2}Q_{M,j}\bigl(z_{j}\mid\operatorname{mask}_{\{j\}}(w(z))\bigr),
\qquad
\hat z=\hat z^{(m^{\star})},\quad
m^{\star}=\min\Bigl(\operatorname*{arg\,max}_{0\le m<R}\Lambda\bigl(\hat z^{(m)}\bigr)\Bigr).
\end{equation}
%
A sweep is one pass over $\mathcal{S}$ in increasing order, a restart is settled when a whole pass
changes nothing, and the budget is six passes for \eqref{eq:remask} and four for the two
vocabularies below. A restart that has not settled when the budget runs out is \emph{not}
discarded: its last state is scored like any other, so the budget is part of the reported rate.
Selected-output sweep counts differ sharply across vocabularies. Of $22{,}714$
reported re-masking probe/context outputs, $950$ selected states use all six sweeps;
$95.8\%$ therefore stopped before the cap. These are not counts of every random restart.
A six-sweep count alone does not distinguish convergence on the final pass from
budget exhaustion, so we do not certify that remainder as fixed points. The
historical BPE selected state uses all four passes in each of its $1{,}750$ outputs,
so those records do not establish that its output is a fixed point. This count
belongs to the historical search campaign, not to the separate 60-cell
common-interval correction audit of \eqref{eq:bpespan}. These stopping counts describe the
search outputs; they do not certify a recovery bound under the separately evaluated LOO probe.

Because $\Lambda$ uses only information available to the adversary, the reported rate measures an
attack rather than oracle-assisted selection. We also record the oracle's best-of-$R$ result. The
difference quantifies the cost of selecting an output without knowing the answer. $\Lambda$ is a sum over slots and, alone among the scores here,
is not divided by nucleotides, which is admissible only because $\mathcal{S}$ is identical
across the $R$ restarts being compared.

We adapt \eqref{eq:remask} for two vocabularies that do not admit it as written. For overlapping \kmer{k}s, independently updating a token can violate consistency
with its overlapping neighbors; the latent variable is therefore the nucleotide, and to re-mask position $j$ we
mask the $k$ tokens covering it and marginalize their predictions onto the letter there,
%
\begin{equation}
\label{eq:ntmarg}
P_{M}\bigl(h_{j}=a\mid\cdot\bigr)\propto
\exp\log\!\!\sum_{v\in K_{j,a}}\!\!\exp\eta_{j,v},
\qquad a\in\Sigma,
\end{equation}
%
with $\eta_{j,\cdot}$ the logits at that slot and $K_{j,a}$ the $k$-mers consistent with the
visible letters that spell $a$ at $j$. Under byte-pair encoding, a nucleotide substitution can change the slot grid because it re-segments
the window. We therefore score the
four candidates on a common interval, grown until it is stable under all of them,
%
\noindent See \eqref{eq:bpespan} in the main text.

%
starting from $I^{(0)}_{j}=[j,j+1)$, where $w^{\,j\leftarrow a}$ is the window with its $j$-th
nucleotide overwritten by $a$. The interval $I_{j}$ depends on $j$ and on the window away from $j$,
never on which letter is written there, so all four candidates are scored on the same nucleotides
and $\Psi_{j}$ divides by nothing. Our initial implementation instead divided by the nucleotide span covered by each candidate's own
tokens. This introduces a prior over letters rather than a normalization,
and on this vocabulary it is large enough to decide the output (appendix~\S\ref{app:bpenorm}).


\paragraph{Candidate-score sensitivity.}
The rescaling $2L/n_{y}(o)$ is a shared constant under a
fixed-width tokenizer but a \emph{per-candidate} divisor under byte-pair encoding, where the covered
span ranges over $96$ to $107$ nucleotides across the $101$ candidates of a cell; it is the divisor
appendix~\S\ref{app:bpenorm} shows to be a prior rather than a normalization, acting here on real segments rather
than single letters. The planted records and pool candidates have mean covered spans of $100.2$ and $99.7$, respectively. The never-trained control uses the same scorer, but this alone does not rule out an interaction between the divisor and membership. This remains a defect of the scoring rule, so we quantified its effect directly: from one forward pass per cell we re-ranked every candidate under the divided
score above, the undivided total, and a \emph{span-matched} score that attributes each token's
surprise across the nucleotides it covers and sums over $[o,o+2L)$ alone, scoring every candidate on
exactly the target's nucleotides. Across 2{,}000 record evaluations (five run/checkpoint pairs,
the never-trained controls and the 1-, 4- and 16-copy tiers, both host conditions, fifty records
each), the span-matched score moves the rank-1 verdict in 33 cells
and the undivided total in 20, $1.7\%$ and $1.0\%$; the single-copy fresh-window rate of
\S\ref{sec:ident} reads $0.18$--$0.26$ across seeds under the rule we report and $0.22$--$0.28$
span-matched. At the final chromosome-22 checkpoint, the three seeds give member rates
$18/22/26\%$ under the reported scorer and $22/22/28\%$ span-matched; matched fresh-window
controls give $0/0/0\%$ and $2/0/0\%$, respectively (50 records per cell). Separately, pooling
the 500 control evaluations across all five run/checkpoint pairs and both hosts gives
$0.014$ under the reported scorer, $0.014$ span-matched and $0.016$ without the divisor.
The five pairs are chromosome-22 seeds 0--2 at final and \emph{E.~coli} seed~0 at rounds~30
and~150 (final). Where the rules disagree, the span-matched rule reports \emph{more} leakage.
The tested cells retain positive membership signal under these alternative scores; the numerical rates remain scorer-dependent.

\section{Public input release identifiers}
\label{app:resources}
The IGSR GRCh38 phased call set was announced on 12 March 2019. The chromosome-22
file, \path{ALL.chr22.GRCh38.phased.vcf.gz}, is dated 27 February 2019 and declares
\texttt{source=IGSRpipeline} and reference
\path{GRCh38_full_analysis_set_plus_decoy_hla.fa}. Its header lists the 2{,}548
samples used for cohort counts. The recall is described by
\citet{lowygallego2019grch38}; the phase-3 study is \citet{thousandgenomes2015}.

\section{Source cells for cross-family recovery comparisons}
\label{app:sourcecells}
Tables~\ref{tab:hyena30} and~\ref{tab:overlap50} supply the cells cited in
\S\ref{sec:slots} and the Discussion. The source paths, configurations and file hashes
are in \path{results/appendix_source_cells.json}; \path{pilot2/appendix_source_cells.py}
reproduces the counts from saved outputs without model inference.
For the uniform-background HyenaDNA runs, \texttt{scores\_final.json} records epoch~30;
for the other backgrounds, the matching file is \texttt{scores\_ep30.json}.
Thus Table~\ref{tab:hyena30} compares one actual training epoch, not one shared filename.
The chromosome-22 single-copy HyenaDNA result in \S\ref{sec:ident} is a separate
human-record campaign with 50 records per tier.

% Verified by pilot2.appendix_source_cells; source hashes in results/appendix_source_cells.json.
\begin{table}[tbp]
\centering
\caption{HyenaDNA medium-160k at epoch 30, learning rate $2\times10^{-5}$.
Exact recovery of the hidden 48\nt{} by one greedy prefix continuation.
Canaries are uniform synthetic 96\nt{} DNA; members are planted sixteen times at random offsets.
Member results use the own window, with 20 records
per seed. Counts follow the listed seed order; fresh-window controls are never trained.}
\label{tab:hyena30}
\begin{tabular}{lllc}
\toprule
Background & Seeds & Member hits/evaluated & Control hits/evaluated \\
\midrule
Uniform DNA & 0, 1, 2 & 14/20, 19/20, 13/20 & 0/40 in each seed \\
\emph{E.~coli} & 0, 1 & 17/20, 14/20 & 0/20 in each seed \\
Yeast & 0, 1, 2 & 18/20, 16/20, 20/20 & 0/20 in each seed \\
GUE promoters & 0, 1, 2 & 16/20, 17/20, 13/20 & 0/20 in each seed \\
\bottomrule
\end{tabular}
\end{table}

% Verified by pilot2.appendix_source_cells; source hashes in results/appendix_source_cells.json.
\begin{table}[tbp]
\centering
\caption{Overlapping DNABERT, masked fine-tuning on chromosome 22, epoch 120,
seed 0: one greedy prefix continuation, with exactness evaluated on the hidden 48\nt{}.
Each tokenizer evaluates the same 50 sixteen-copy human-variant-derived records in both
access conditions. These are the full \texttt{scores\_ep120\_n50.json} evaluations.}
\label{tab:overlap50}
\begin{tabular}{ccccc}
\toprule
$k$ & Own hits/evaluated & Own rate & Fresh hits/evaluated & Fresh rate \\
\midrule
3 & 0/50 & 0\% & 0/50 & 0\% \\
4 & 1/50 & 2\% & 1/50 & 2\% \\
5 & 3/50 & 6\% & 1/50 & 2\% \\
6 & 25/50 & 50\% & 18/50 & 36\% \\
\bottomrule
\end{tabular}
\end{table}

